\chapter{Inductive types and iota-reduction}
\label{chap:inductive}

This chapter covers the model of Lean's inductive declarations: the syntactic shapes the
constants of an inductive block must have, the staged environment extension that installs
them, the well-formedness record they are checked against, and the machinery that turns a
recursor computation rule into a registered schematic reduction rule ($\iota$ reduction).
Four files are in scope: \srcloc{Lean4Lean/Theory/Inductive.lean}{1}, the specification
layer (444 lines); \srcloc{Lean4Lean/Theory/Typing/Pattern.lean}{1}, the abstract pattern
machinery together with the $\iota$ reduct builder (725 lines);
\srcloc{Lean4Lean/Theory/Typing/InductiveLemmas.lean}{1}, the environment bookkeeping and
the soundness theorem for \texttt{addInduct} (785 lines); and
\srcloc{Lean4Lean/Theory/Typing/InductiveParams.lean}{1}, a concrete instance of the
Church-Rosser development's parameter record (433 lines).

The baseline matters for judging this material. On \texttt{master},
\texttt{Theory/Inductive.lean} was seven lines, consisting of
\texttt{def VInductDecl.WF := sorry} and \texttt{def VEnv.addInduct := sorry}, and
\texttt{InductiveLemmas.lean} was a single \texttt{theorem addInduct\_WF ... := sorry}.
Inductive types were entirely unmodelled: \texttt{VInductDecl.WF} was an opaque proposition,
so no environment containing an inductive declaration could be shown well-formed and every
downstream theorem quantifying over well-formed environments was uninstantiable there. The
\texttt{iota} branch
replaces both \texttt{sorry}-definitions with real definitions, proves
\texttt{addInduct\_WF} relative to the new record, makes the $\iota$ rule a first-class
reduction rule of the type system, and builds the first concrete instance of the
\texttt{Params} class the Church-Rosser development is parameterised over. Counting by
line attribution: \texttt{Inductive.lean} is 4 master / 438 \texttt{iota} / 2
\texttt{trproj}; \texttt{InductiveLemmas.lean} 8 / 753 / 24;
\texttt{InductiveParams.lean} 0 / 389 / 44; \texttt{Pattern.lean} 229 / 419 / 77. In total
about 2150 of the group's 2387 lines are contributed, and 69 of this chapter's 74 nodes.

Thesis correspondence is \texttt{axioms.tex} \S2.6: \S2.6.1 inductive specifications
(\texttt{sec:inductive}), \S2.6.2 large elimination (\texttt{sec:large\_elim}), \S2.6.3 the
recursor, \S2.6.4 the computation rule (\texttt{sec:iota}). The model deviates from the
thesis in two systematic, deliberate, kernel-facing directions. It carries a
parameter/index split (\texttt{decl.nparams}) that the thesis's $t\;p[b]$ does not have,
and it supports mutual blocks (\texttt{decl.types} a list) where the thesis has a single
$\mu t{:}F.\,K$. Both are recorded in the docstrings. A third and more consequential
difference is that the thesis builds positivity into the grammar of a \texttt{ctor}
judgment, whereas the model states flat predicates over an already elaborated $\Pi$-type,
mirroring the kernel's \texttt{checkPositivity} and \texttt{isValidIndApp?}; no lemma
anywhere relates the two.

An attribution caveat applies throughout. On these four files
\texttt{git diff iota trproj} consists of two line re-wraps
(\srcloc{Lean4Lean/Theory/Inductive.lean}{397},
\srcloc{Lean4Lean/Theory/Typing/InductiveLemmas.lean}{618}) and two docstring rewrites
(\srcloc{Lean4Lean/Theory/Typing/InductiveParams.lean}{351},
\srcloc{Lean4Lean/Theory/Typing/InductiveParams.lean}{385}): every \emph{code} block the
line blame tags \texttt{trproj} was ported to \texttt{iota} as commit \texttt{3de1dcb},
byte for byte. The blame's tag is kept in the badges below, but for review purposes the
code of this group belongs to the $\iota$ contribution. What \texttt{trproj} has here that
\texttt{iota} does not is the expanded \texttt{DefEqsAsPats}/\texttt{toParams}
documentation, which is the most useful prose in the group
(\ref{def:indparams-defeqs-as-pats}).

Assessment. The design is good and unusually well considered: the staging of
\texttt{VInductDecl.WF} so that each kind of constant is typed in the environment the
kernel checks it in, the \texttt{Generic}/\texttt{Realizes}/\texttt{TemplateHeaded} triple
that makes the new \texttt{IsDefEq.pat} rule legal and informative, the split of
\texttt{iotaRHS'} (what reduction sees) from \texttt{iotaRHS} (the typed split), and the
\texttt{PatsIota} registry invariant are all the right abstractions, each justified in
prose in the source. No file in the group contains a \texttt{sorry}, the proofs are short,
and the specification is exercised against real kernel data by \texttt{Tests/IotaShape.lean}
through the \texttt{Decidable} instances of \texttt{Tests/ShapeDecide.lean}. Three things
limit what the result means, all recorded in \S\ref{sec:ind-review}: nothing derives
\texttt{VInductDecl.WF} from the kernel's own inductive checker, so the hardest declaration
kind remains an assumption of the refinement; \texttt{VExpr.RuleShape} pins only the
\emph{number} of recursive arguments of an $\iota$ reduct, so the record admits rules Lean
would never generate; and the \texttt{Params} instance applies only to environments with no
definitional axiom at all. The contribution should be described as specified and validated,
not as verified.

\section{Syntactic shapes of an inductive block}

All predicates in this section live in namespace \texttt{VExpr} and read an already
elaborated type. Notation: \texttt{bvarsDesc lo n} is the descending variable list
$[\mathsf{bvar}(lo+n-1),\ldots,\mathsf{bvar}\,lo]$, \texttt{piArity}/\texttt{piBinders}/\texttt{piBody}
read a $\Pi$-telescope and \texttt{lamArity}/\texttt{lamBody} a $\lambda$-telescope.

\begin{definition}[Occurrence of a block's type formers, \texttt{MentionsConst}]
\label{def:ind-mentions-const}
\lean{Lean4Lean.VExpr.MentionsConst, Lean4Lean.VExpr.mentionsConst, Lean4Lean.VExpr.mentionsConst_iff}
\leanok
\contrib{iota} \srcloc{Lean4Lean/Theory/Inductive.lean}{24}
\uses{ind:vexpr}
$\mathsf{MentionsConst}\;cs\;e$ is defined by recursion on $e$: false at $\mathsf{bvar}$
and $\mathsf{sort}$, list membership at $\mathsf{const}$, disjunction at the three binary
nodes. \texttt{mentionsConst} is the Boolean twin and \texttt{mentionsConst\_iff} relates
them by a structural induction discharged by \texttt{simp}.
\thesisref{axioms.tex \S2.6.1 (\texttt{sec:inductive}), the implicit side condition that a
nonrecursive argument type does not mention $t$}
\reviewnote{It mirrors the kernel's \texttt{hasIndOcc}
(\srcloc{Lean4Lean/Inductive/Add.lean}{118}), which searches the whole term, so the two
agree. The Boolean copy exists only so that \texttt{Tests/ShapeDecide.lean} can
\texttt{decide} the shape predicates; the \texttt{Decidable} instance built from
\texttt{mentionsConst\_iff} lives there, not beside the definition, a placement the test
file's header states and defends.}
\end{definition}

\begin{definition}[Result type of a constructor, \texttt{CtorResult}]
\label{def:ind-ctor-result}
\lean{Lean4Lean.VExpr.CtorResult}
\leanok
\contrib{iota} \srcloc{Lean4Lean/Theory/Inductive.lean}{44}
\uses{def:pi-telescope, def:bvarsdesc, def:mkapps}
$ty.\mathsf{CtorResult}\;T\;np\;nf\;nind$ holds when $ty.\mathsf{piArity} = np+nf$ and there
are universe levels $us$ and a list $idx$ of length $nind$ with
$ty.\mathsf{piBody} = (\mathsf{const}\;T\;us).\mathsf{mkApps}\,(\texttt{bvarsDesc}\;nf\;np \mathbin{+\!\!+} idx)$.
Under $np+nf$ binders the parameters sit at $\mathsf{bvar}(nf+np-1),\ldots,\mathsf{bvar}\,nf$,
which is exactly \texttt{bvarsDesc nf np}; the index terms are unconstrained here.
\thesisref{axioms.tex \S2.6.1 (\texttt{sec:inductive}), the base case $t\;e\;\mathsf{ctor}$
of the \texttt{ctor} judgment}
\reviewnote{The thesis has one argument sequence where the model fixes the first
\texttt{nparams} arguments to be the parameter variables. This is a deliberate move towards
the kernel and is documented as such.}
\end{definition}

\begin{definition}[A valid recursive occurrence of the block, \texttt{ValidIndApp}]
\label{def:ind-valid-ind-app}
\lean{Lean4Lean.VExpr.ValidIndApp}
\leanok
\contrib{iota} \srcloc{Lean4Lean/Theory/Inductive.lean}{75}
\uses{def:ind-mentions-const, def:bvarsdesc, def:mkapps}
$e.\mathsf{ValidIndApp}\;fs\;np\;d$ holds when there is a former $T \in fs$, levels $us$ and
index terms $idx$ with
$e = (\mathsf{const}\;T\;us).\mathsf{mkApps}\,(\texttt{bvarsDesc}\;d\;np \mathbin{+\!\!+} idx)$
and no $a \in idx$ mentions a former of $fs$. The offset $d$ is the number of binders
between the field level and the occurrence.
\thesisref{axioms.tex \S2.6.1 (\texttt{sec:inductive}), the $t\;e$ of the recursive-argument
rule}
\reviewnote{The kernel's \texttt{isValidIndAppIdx} (\srcloc{Lean4Lean/Inductive/Add.lean}{157})
additionally checks the exact arity \texttt{args.size == nparams + nindices[i]}. The model
drops it and so accepts a former over-applied in its index positions. Such a field would be
rejected later by \texttt{WF.universes}, so this is imprecision rather than a hole, but the
docstring advertises the predicate as mirroring \texttt{isValidIndApp?} and it does not.}
\end{definition}

\begin{definition}[Strict positivity of one constructor field, \texttt{FieldPositive}]
\label{def:ind-field-positive}
\lean{Lean4Lean.VExpr.FieldPositive}
\leanok
\contrib{iota} \srcloc{Lean4Lean/Theory/Inductive.lean}{93}
\uses{def:ind-valid-ind-app, def:ind-mentions-const, def:pi-telescope}
A field type $ty$ is positive at depth $d$ when either it does not mention the block at
all, or every one of its $\Pi$-binder types is free of the block's formers and
$ty.\mathsf{piBody}$ is a \texttt{ValidIndApp} at depth $d + ty.\mathsf{piArity}$.
\thesisref{axioms.tex \S2.6.1 (\texttt{sec:inductive}), the two argument rules
$\forall y{:}\beta.\,\tau$ and $(\forall z{::}\gamma.\;t\;e)\to\tau$}
\reviewnote{Two deviations. The kernel reduces the field type and each binder domain to
weak head normal form (\texttt{checkPositivity},
\srcloc{Lean4Lean/Inductive/Add.lean}{184}) whereas the model reads manifest binders only,
so the model is strictly \emph{stricter} than the kernel. Conversely the thesis forbids a
recursive argument from being referred to by later arguments and the model, like the
kernel, does not.}
\end{definition}

\begin{definition}[Strict positivity of a constructor, \texttt{CtorPositive}]
\label{def:ind-ctor-positive}
\lean{Lean4Lean.VExpr.CtorPositive}
\leanok
\contrib{iota} \srcloc{Lean4Lean/Theory/Inductive.lean}{101}
\uses{def:ind-field-positive, def:ind-mentions-const, def:pi-telescope}
$ty.\mathsf{CtorPositive}\;fs\;np$ holds when no former of $fs$ occurs in the first $np$
$\Pi$-binder types, and for every $i < ty.\mathsf{piArity} - np$ the binder at position
$np+i$ exists and is \texttt{FieldPositive} at depth $i$. The depth bookkeeping is correct:
under field $i$'s own $k$ binders the parameters sit at \texttt{bvarsDesc (i + k) np}, which
is what \texttt{FieldPositive} passes on as $d + ty.\mathsf{piArity}$.
\thesisref{axioms.tex \S2.6.1 (\texttt{sec:inductive})}
\end{definition}

\begin{definition}[A field appears among the result indices, \texttt{FieldInIndices}]
\label{def:ind-field-in-indices}
\lean{Lean4Lean.VExpr.FieldInIndices}
\leanok
\contrib{iota} \srcloc{Lean4Lean/Theory/Inductive.lean}{107}
\uses{def:app-spine, def:pi-telescope}
$ty.\mathsf{FieldInIndices}\;np\;i$ says
$\mathsf{bvar}\,(ty.\mathsf{piArity} - np - 1 - i)$ occurs in
$ty.\mathsf{piBody}.\mathsf{getAppArgs}.\mathsf{drop}\;np$.
\thesisref{axioms.tex \S2.6.2 (\texttt{sec:large\_elim}), the side condition $y \in e$}
\reviewnote{The kernel's \texttt{isLargeEliminator}
(\srcloc{Lean4Lean/Inductive/Add.lean}{275}) searches all result arguments, not only those
past \texttt{nparams}. The two agree because \texttt{CtorResult} pins the first $np$
arguments to be \texttt{bvarsDesc nf np}, whose entries are all $\ge nf$, while the field
variable sought here is $\mathsf{bvar}\,(nf-1-i) < nf$; that argument is what makes the
\texttt{drop np} sound and it is recorded nowhere in the source.}
\end{definition}

\begin{definition}[The de Bruijn context of a constructor field, \texttt{fieldCtx}]
\label{def:ind-field-ctx}
\lean{Lean4Lean.VExpr.fieldCtx}
\leanok
\contrib{iota} \srcloc{Lean4Lean/Theory/Inductive.lean}{112}
\uses{def:pi-telescope}
$ty.\mathsf{fieldCtx}\;np\;i = (ty.\mathsf{piBinders}.\mathsf{take}\,(np+i)).\mathsf{reverse}$,
the parameter and earlier-field binder types innermost first, which is the orientation
contexts have in this development. Used only by \texttt{LargeElim} and
\texttt{WF.universes}.
\thesisref{no direct counterpart; a helper for the typing clauses of \S2.6.2}
\end{definition}

\begin{definition}[Type of a recursor's major premise, \texttt{MajorApp}]
\label{def:ind-major-app}
\lean{Lean4Lean.VExpr.MajorApp}
\leanok
\contrib{iota} \srcloc{Lean4Lean/Theory/Inductive.lean}{63}
\uses{def:bvarsdesc, def:mkapps}
$A.\mathsf{MajorApp}\;T\;np\;nm\;nmin\;nind$ says
$A = (\mathsf{const}\;T\;us).\mathsf{mkApps}\,(\texttt{bvarsDesc}\,(nm{+}nmin{+}nind)\;np \mathbin{+\!\!+} \texttt{bvarsDesc}\;0\;nind)$
for some $us$. The arithmetic checks out: at the major binder there are $np+nm+nmin+nind$
binders in scope.
\thesisref{axioms.tex \S2.6.3, the binder $z : P\;a$ of $\mathrm{rec}_P$}
\end{definition}

\begin{definition}[Shape of a motive binder, \texttt{MotiveShape}]
\label{def:ind-motive-shape}
\lean{Lean4Lean.VExpr.MotiveShape}
\leanok
\contrib{iota} \srcloc{Lean4Lean/Theory/Inductive.lean}{121}
\uses{def:headconst, def:pi-telescope}
A motive type is a $\Pi$-telescope whose body is a sort and whose \emph{last} binder is
headed by a constant, that constant being \texttt{motiveFormer?}. Which constant is not
pinned here, because a bare \texttt{VEnv} has no notion of type former; \texttt{RecShape}
ties only the eliminated motive's head to the major premise.
\thesisref{axioms.tex \S2.6.3, $\kappa = \forall a{::}\alpha.\;P\;a \to \mathcal{U}_u$}
\reviewnote{A deliberately weak approximation of $\kappa$: the telescope is not required to
be $indices \to P\;a \to \mathsf{Sort}$. Its ``ends in a sort'' clause is subsumed by
\texttt{VInductDecl.WF.recs\_elim}, which pins the exact sort, so it is also mildly
redundant.}
\end{definition}

\begin{definition}[A minor premise targets one of the motives, \texttt{MinorHeaded}]
\label{def:ind-minor-headed}
\lean{Lean4Lean.VExpr.MinorHeaded}
\leanok
\contrib{iota} \srcloc{Lean4Lean/Theory/Inductive.lean}{128}
\uses{def:app-spine, def:pi-telescope}
$A.\mathsf{MinorHeaded}\;i\;nm$ says $A.\mathsf{piBody}.\mathsf{getAppFn} =
\mathsf{bvar}\,(A.\mathsf{piArity}+i+k)$ for some $k < nm$, that is, the body is headed by
motive $nm-1-k$. The index computation is correct: at minor $i$ there are $np+nm+i$ binders
in scope outside the minor's own telescope.
\thesisref{axioms.tex \S2.6.3, the head $C$ of $\varepsilon_c = \forall b{::}\beta.\forall v{::}\delta.\;C\;p[b]\;(c\;b)$}
\reviewnote{Nothing ties minor $i$ to the motive of the type former its constructor belongs
to, only to \emph{some} motive.}
\end{definition}

\begin{definition}[The minor premise belonging to a constructor, \texttt{MinorFor}]
\label{def:ind-minor-for}
\lean{Lean4Lean.VExpr.MinorFor}
\leanok
\contrib{iota} \srcloc{Lean4Lean/Theory/Inductive.lean}{134}
\uses{def:recheaded, def:headconst, def:app-spine}
$A.\mathsf{MinorFor}\;c$ says $A$ is \texttt{RecHeaded} (its $\Pi$-body's head is a bound
variable) and the \emph{last} argument of that application is headed by the constant $c$.
Its only use in the theory is inside \texttt{VInductDecl.WF.rule\_shape}, which picks, for
each rule, a minor satisfying it.
\thesisref{axioms.tex \S2.6.3, the argument $c\;b$ of $\varepsilon_c$}
\reviewnote{Two distinct constructors cannot satisfy \texttt{MinorFor} for the same minor
(\texttt{getLast?} and \texttt{headConst?} are functions), so the constructor-to-minor map
\texttt{rule\_shape} chooses is injective, but no lemma in the repository derives that and
nothing uses it; \texttt{rules\_nodup} is what the proofs actually appeal to.}
\end{definition}

\begin{definition}[Shape of a recursor type, \texttt{RecShape}]
\label{def:ind-rec-shape}
\lean{Lean4Lean.VExpr.RecShape}
\leanok
\contrib{iota} \srcloc{Lean4Lean/Theory/Inductive.lean}{143}
\uses{def:ind-motive-shape, def:ind-minor-headed, def:ind-major-app, def:pi-telescope, def:mkapps, def:bvarsdesc}
$ty.\mathsf{RecShape}\;np\;nm\;nmin\;nind$ requires $ty.\mathsf{piArity} = np+nm+nmin+nind+1$;
the $nm$ binders after the parameters to be \texttt{MotiveShape}; the next $nmin$ to be
\texttt{MinorHeaded}; and, for some $j < nm$, the last binder to be $T$ applied to the
parameter and index variables (\texttt{MajorApp}) with $T$ the head of motive $j$'s last
binder, while $ty.\mathsf{piBody}$ is $\mathsf{bvar}\,(nind+nmin+1+(nm-1-j))$ applied to
\texttt{bvarsDesc 0 (nind+1)}, that is, motive $j$ applied to the indices and the major.
\thesisref{axioms.tex \S2.6.3, $\mathrm{rec}_P : \forall C{:}\kappa.\forall e{::}\varepsilon.\forall a{::}\alpha.\forall z{:}P\;a.\;C\;a\;z$}
\reviewnote{Parameter and index binder types are unconstrained, as are the heads of the
non-eliminated motives of a mutual recursor: a documented weakening of $\kappa$ and
$\varepsilon$. The de Bruijn arithmetic of the body head is machine-checked against the
kernel by \texttt{checkShapes} in \texttt{Tests/IotaShape.lean}, which decides
\texttt{RecShape} on the translated \texttt{RecursorVal.type} of some fifty recursors
(\texttt{Nat}, \texttt{Eq}, \texttt{Acc}, \texttt{Vec}, the \texttt{Ev}/\texttt{Od} mutual
block, \texttt{False}, and a list of \texttt{Init}/\texttt{Std} inductives).}
\end{definition}

\begin{definition}[Shape of a constructor type, \texttt{CtorShape}]
\label{def:ind-ctor-shape}
\lean{Lean4Lean.VExpr.CtorShape}
\leanok
\contrib{iota} \srcloc{Lean4Lean/Theory/Inductive.lean}{157}
\uses{def:ctorheaded, def:pi-telescope}
$ty.\mathsf{CtorShape}\;arity$ is $ty.\mathsf{piArity} = arity$ together with
$ty.\mathsf{CtorHeaded}$. It is deliberately weaker than \texttt{CtorResult}: it names
neither the type former nor the index count, which is what lets
\ref{thm:indlem-rules-ctor-shape} state it about a constant looked up in an environment
where the block is no longer in hand.
\thesisref{no direct counterpart; the environment-level residue of \S2.6.1}
\reviewnote{Its two halves are consumed separately and neither by the same client:
\texttt{PatsIota.ctor\_shape} keeps only the \texttt{CtorHeaded} half, and the front end
(\srcloc{Lean4Lean/Verify/Environment/Basic.lean}{514}) only the arity half.}
\end{definition}

\begin{definition}[The type former a recursor eliminates, \texttt{majorFormer?}]
\label{def:ind-major-former}
\lean{Lean4Lean.VExpr.majorFormer?}
\leanok
\contrib{iota} \srcloc{Lean4Lean/Theory/Inductive.lean}{161}
\uses{def:headconst, def:pi-telescope, def:getmajoridx}
$ty.\mathsf{majorFormer?}\;idx = ty.\mathsf{piBinders}[idx]?.\mathsf{bind}\;\mathsf{headConst?}$,
the head constant of the $idx$-th $\Pi$-binder, that is, of the major premise when
$idx = r.\mathsf{getMajorIdx}$.
\thesisref{axioms.tex \S2.6.3, the $t$ eliminated by $\mathrm{rec}_P$}
\end{definition}

\begin{definition}[Shape of an iota-rule reduct, \texttt{RuleShape}]
\label{def:ind-rule-shape}
\lean{Lean4Lean.VExpr.RuleShape}
\leanok
\contrib{iota} \srcloc{Lean4Lean/Theory/Inductive.lean}{174}
\uses{def:lam-telescope, def:bvarsdesc, def:mkapps}
$rhs.\mathsf{RuleShape}\;np\;nm\;nmin\;nf\;nrec\;j$ says
$rhs.\mathsf{lamArity} = np+nm+nmin+nf$ and, for some list $recArgs$ of length $nrec$,
$rhs.\mathsf{lamBody} = \mathsf{bvar}\,(nf + (nmin-1-j))$ applied to
$\texttt{bvarsDesc}\;0\;nf \mathbin{+\!\!+} recArgs$: minor $j$ applied $\eta$-long to the
fields and then to $nrec$ further arguments.
\thesisref{axioms.tex \S2.6.4 (\texttt{sec:iota}), $\mathrm{rec}_P\;C\;e\;p[b]\;(c\;b) \equiv e_c\;b\;v$}
\reviewnote{This is the main weakening of \S2.6.4 in the file, flagged at length in its own
docstring: the recursive calls $v_i = \lambda x{::}\xi_i.\;\mathrm{rec}_P\;C\;e\;\pi_i[b,x]\;(u_i\;x)$
are pinned only in number, not in shape, and \texttt{WF.rules\_wf} only forces them to be
well-typed. \texttt{VInductDecl.WF} therefore accepts $\iota$ rules Lean would never
generate.}
\end{definition}

\begin{lemma}[Spine reformulations of the shape predicates]
\label{family:ind-shape-iff}
\lean{Lean4Lean.VExpr.CtorResult_iff, Lean4Lean.VExpr.MajorApp_iff, Lean4Lean.VExpr.ValidIndApp_iff}
\leanok
\contrib{iota} \srcloc{Lean4Lean/Theory/Inductive.lean}{48}
\uses{def:ind-ctor-result, def:ind-major-app, def:ind-valid-ind-app}
Each of \texttt{CtorResult}, \texttt{MajorApp} and \texttt{ValidIndApp} is re-expressed
through \texttt{headConst?}, \texttt{getAppArgs} and a list-prefix condition, eliminating
the existentials over the universe list and the index list.
\thesisref{no direct counterpart; a decidability interface for \S2.6.1 and \S2.6.3}
\end{lemma}
\begin{proof}
\uses{fam:const-spine, def:bvarsdesc}
\leanok
All three are rewrites by \texttt{eq\_const\_mkApps\_iff} and
\texttt{eq\_const\_mkApps\_append\_iff} together with \texttt{bvarsDesc\_length}; the
\texttt{CtorResult} case first commutes the two conjuncts under the existentials.
Their only consumer is \texttt{Tests/ShapeDecide.lean}, where they supply the
\texttt{Decidable} instances, so the theory files carry three lemmas whose sole purpose is
testing.
\end{proof}

\begin{lemma}[Immediate consequences of the shape predicates]
\label{family:ind-shape-corollaries}
\lean{Lean4Lean.VExpr.RuleShape.lam, Lean4Lean.VExpr.RecShape.recHeaded, Lean4Lean.VExpr.RecShape.one_le_numMotives, Lean4Lean.VExpr.RecShape.majorFormer?_eq, Lean4Lean.VExpr.CtorResult.ctorHeaded, Lean4Lean.VExpr.CtorResult.ctorShape}
\leanok
\contrib{iota} \srcloc{Lean4Lean/Theory/Inductive.lean}{183}
\uses{def:ind-rule-shape, def:ind-rec-shape, def:ind-ctor-result, def:ind-ctor-shape, def:ind-major-former, def:recheaded, def:ctorheaded}
A \texttt{RuleShape} reduct is a $\lambda$-abstraction as soon as $j < nmin$; a
\texttt{RecShape} type is \texttt{RecHeaded}, has at least one motive, and its major
premise's head constant is the one \texttt{majorFormer?} reads; a \texttt{CtorResult} type
is \texttt{CtorHeaded} and has \texttt{CtorShape (np + nf)}.
\thesisref{no direct counterpart; interface lemmas}
\end{lemma}
\begin{proof}
\leanok
Each is a projection out of the defining conjunction or existential, with \texttt{omega} for
the arity case of \texttt{RuleShape.lam} and \texttt{getAppFn\_mkApps} for the two
head-shape cases.
\reviewnote{\texttt{RuleShape.lam} (used by \texttt{addRules\_ordered}),
\texttt{RecShape.recHeaded} (used by \texttt{PatsIota.induct}) and
\texttt{CtorResult.ctorShape} (used by \texttt{rules\_ctor\_shape}) are load-bearing.
\texttt{RecShape.one\_le\_numMotives} (\srcloc{Lean4Lean/Theory/Inductive.lean}{196}) and
\texttt{RecShape.majorFormer?\_eq} (\srcloc{Lean4Lean/Theory/Inductive.lean}{202}) are
referenced nowhere in the repository.}
\end{proof}

\section{Large elimination}

\begin{definition}[Large elimination, \texttt{LargeElim}]
\label{def:ind-large-elim}
\lean{Lean4Lean.VInductDecl.LargeElim}
\leanok
\contrib{iota} \srcloc{Lean4Lean/Theory/Inductive.lean}{226}
\uses{def:ind-field-in-indices, def:ind-field-ctx, fam:vlevel-lattice, def:hastype-istype, struct:vinductdecl}
$decl.\mathsf{LargeElim}\;env\;\ell$ is a three-way disjunction: $\ell$ is never zero; or the
block is a single type former with no constructor; or a single type former with one
constructor each of whose fields is either typed at $\mathsf{sort}\;0$ in $env$ and the
field's own \texttt{fieldCtx}, or occurs among the indices of the constructor's result type.
Here $env$ is the stage-0 environment, the one in which the type formers are declared.
\thesisref{axioms.tex \S2.6.2 (\texttt{sec:large\_elim}), the $\mathrm{LE}$ judgments}
\reviewnote{It matches the kernel's \texttt{isLargeEliminator}
(\srcloc{Lean4Lean/Inductive/Add.lean}{258}) clause for clause, including the zero-constructor
case, and the thesis's separate rule for recursive arguments is subsumed. But the kernel
tests \texttt{isAlwaysZero} on the inferred level where the model demands a typing at
$\mathsf{sort}\;0$; these agree only up to level defeq and uniqueness of typing, neither of
which is proved here. Because the propositionality clause is a typing judgment,
\texttt{LargeElim} is undecidable.}
\end{definition}

\begin{definition}[The decidable syntactic half of large elimination, \texttt{LargeElimShape}]
\label{def:ind-large-elim-shape}
\lean{Lean4Lean.VInductDecl.LargeElimShape, Lean4Lean.VInductDecl.LargeElim.shape}
\leanok
\contrib{iota} \srcloc{Lean4Lean/Theory/Inductive.lean}{237}
\uses{def:ind-large-elim, struct:vinductdecl}
$decl.\mathsf{LargeElimShape}$ says the block has exactly one type former and every type
former has at most one constructor. \texttt{LargeElim.shape} proves that a block which
eliminates largely while its result sort may be $\mathsf{Prop}$ has that shape, by a
three-way case split reading the singleton lists off the second and third disjuncts.
\thesisref{axioms.tex \S2.6.2 (\texttt{sec:large\_elim})}
\reviewnote{\texttt{LargeElimShape} carries a \texttt{Decidable} instance and is used by the
tests; \texttt{LargeElim.shape}, the lemma that justifies calling it the syntactic half, is
used nowhere.}
\end{definition}

\section{Patterns, matching and schematic rule reducts}

This section is \srcloc{Lean4Lean/Theory/Typing/Pattern.lean}{1}. The first five nodes are
master's abstract machinery for schematic rewrite rules; the rest is what the $\iota$ rule
needed added to it. The additions are strictly additive: \texttt{git diff master} on this
file shows no deleted lines.

\begin{definition}[Patterns, subpatterns and pattern intersection]
\label{inductive:pattern-core}
\lean{Lean4Lean.Pattern, Lean4Lean.Pattern.varN, Lean4Lean.Subpattern, Lean4Lean.Subpattern.trans, Lean4Lean.Subpattern.sizeOf_le, Lean4Lean.Subpattern.antisymm, Lean4Lean.Arity, Lean4Lean.Arity.subpattern, Lean4Lean.Pattern.inter, Lean4Lean.Pattern.inter_self, Lean4Lean.Pattern.inter_comm, Lean4Lean.Pattern.LE}
\leanok
\srcloc{Lean4Lean/Theory/Typing/Pattern.lean}{8}
\uses{ind:vexpr}
A \texttt{Pattern} is a constant, an application of two patterns, or a \texttt{var} hole over
a pattern; \texttt{varN} iterates \texttt{var}. \texttt{Subpattern} is reflexive structural
containment, \texttt{Arity} counts argument positions above a subpattern, and
\texttt{Pattern.inter} computes the common refinement of two patterns: at an
\texttt{app}/\texttt{var} pair the \texttt{app}'s argument pattern survives
($\mathsf{var}\,f \sqcap \mathsf{app}\,f'\,a' = \mathsf{app}\,(f \sqcap f')\,a'$), since a
\texttt{var} hole matches any argument. Antisymmetry of \texttt{Subpattern} goes through a
\texttt{sizeOf} bound.
\reviewnote{\texttt{Pattern.LE} is declared and used nowhere in the repository, and
\texttt{Arity}/\texttt{Arity.subpattern} only in the experimental logical-relations
development (\srcloc{Lean4Lean/Experimental/ShapeLogRel.lean}{3436} and
\srcloc{Lean4Lean/Experimental/ShapeLogRel.lean}{3620}), never in the verified one. This is
pre-existing \texttt{master} code, not the contributor's.}
\end{definition}

\begin{definition}[Matching a pattern against an expression]
\label{inductive:pattern-matches}
\lean{Lean4Lean.Pattern.Path, Lean4Lean.Pattern.Matches, Lean4Lean.Pattern.Matches.uniq, Lean4Lean.Pattern.matches_determ, Lean4Lean.Pattern.matches_inter, Lean4Lean.Pattern.OnArgs}
\leanok
\srcloc{Lean4Lean/Theory/Typing/Pattern.lean}{132}
\uses{inductive:pattern-core}
$p.\mathsf{Path}$ enumerates the holes of $p$, and $p.\mathsf{Matches}\;e\;ls\;g$ says $e$
matches $p$ with the head constant instantiated at levels $ls$ and hole assignment $g$.
Matching is deterministic (both by structural induction on the derivation), and an
expression matches $p$ and $q$ iff it matches their intersection.
\reviewnote{In the \texttt{app} case the argument side's level list is discarded, so an
$\iota$ redex constrains the recursor's universes but not the constructor's. That is what
makes \texttt{iotaRHS}'s use of the recursor's levels correct.
\texttt{Matches.uniq} (line 143) and \texttt{matches\_determ} (line 283) are the same
statement proved twice.}
\end{definition}

\begin{definition}[Rule reducts and side conditions]
\label{inductive:pattern-rhs}
\lean{Lean4Lean.Pattern.RHS, Lean4Lean.Pattern.Check, Lean4Lean.Pattern.RHS.apply, Lean4Lean.Pattern.Check.OK, Lean4Lean.Pattern.Check.OK.map}
\leanok
\srcloc{Lean4Lean/Theory/Typing/Pattern.lean}{155}
\uses{inductive:pattern-core, inductive:pattern-matches}
A $p.\mathsf{RHS}$ is a tree of closed fixed terms, applications, and references to holes of
$p$; \texttt{RHS.apply} instantiates it at a match, level-instantiating the fixed parts. A
$p.\mathsf{Check}$ is a list of pairs of reducts to be related, and \texttt{Check.OK}
asserts them against a given relation; \texttt{OK.map} transports it along a relation
morphism.
\end{definition}

\begin{lemma}[Lifting and instantiation of matches and reducts]
\label{family:pattern-transport-master}
\lean{Lean4Lean.Pattern.RHS.lift'_apply, Lean4Lean.Pattern.RHS.liftN_apply, Lean4Lean.Pattern.matches_lift', Lean4Lean.Pattern.matches_liftN, Lean4Lean.Pattern.RHS.instN_apply, Lean4Lean.Pattern.matches_instN}
\leanok
\srcloc{Lean4Lean/Theory/Typing/Pattern.lean}{186}
\uses{inductive:pattern-matches, inductive:pattern-rhs}
\texttt{RHS.apply} commutes with lifting and with instantiation (the fixed parts being
closed). For lifting, matching is characterised both ways: $p$ matches $e.\mathsf{lift}'\rho$
iff it matches $e$ with correspondingly lifted holes. For instantiation only the forward
direction is proved: a match of $e$ yields a match of $e.\mathsf{inst}\;e_0\;k$ with
instantiated holes.
\end{lemma}
\begin{proof}
\leanok
Structural induction on the reduct for the \texttt{apply} lemmas; for
\texttt{matches\_lift'} an induction on the \texttt{Matches} derivation generalised over the
pre-image in the forward direction and over the lifted holes in the backward one, with
\texttt{matches\_liftN} a corollary through \texttt{lift'\_consN\_skipN}.
\end{proof}

\begin{definition}[The two schematic rule shapes, \texttt{SimplePattern}]
\label{inductive:simple-pattern}
\lean{Lean4Lean.SimplePattern, Lean4Lean.SimplePattern.toPattern}
\leanok
\srcloc{Lean4Lean/Theory/Typing/Pattern.lean}{484}
\uses{inductive:pattern-core}
A \texttt{SimplePattern} is either $\mathsf{iota}\;rec\;M\;ctor\;N$, unfolding to
$(\mathsf{const}\;rec).\mathsf{varN}\;M$ applied to
$(\mathsf{const}\;ctor).\mathsf{varN}\;N$, or $\mathsf{defn}\;head$, the bare constant.
\reviewnote{\texttt{SimplePattern.defn} is never registered by anything;
\texttt{InductiveParams} explains why, namely that $\delta$ rules stay in \texttt{defeqs}.}
\end{definition}

\begin{lemma}[Combinatorics of constructor spines]
\label{family:pattern-spine-combinatorics}
\lean{Lean4Lean.Pattern.subpattern_varN_const, Lean4Lean.Pattern.not_app_subpattern_varN_const, Lean4Lean.Pattern.varN_const_inter, Lean4Lean.Pattern.varN_const_inj}
\leanok
\contrib{iota} \srcloc{Lean4Lean/Theory/Typing/Pattern.lean}{76}
\uses{inductive:pattern-core}
Every subpattern of $(\mathsf{const}\;c).\mathsf{varN}\;k$ is
$(\mathsf{const}\;c).\mathsf{varN}\;i$ for some $i \le k$; a spine contains no application
subpattern; two spines intersect to \texttt{some} only when head and length agree, and are
equal only then.
\thesisref{no direct counterpart; the combinatorial content of \S2.6.4 as a rewrite system}
\end{lemma}
\begin{proof}
\leanok
Induction on $k$ for the subpattern lemma, from which the no-application case follows by
inversion; \texttt{varN\_const\_inter} is a four-way induction on the two lengths, and
\texttt{varN\_const\_inj} deduces injectivity from it through \texttt{inter\_self}. These
four lemmas are exactly what the \texttt{Params} coherence conditions need later.
\end{proof}

\begin{definition}[The holes a reduct mentions, \texttt{RHS.Uses}]
\label{def:pattern-rhs-uses}
\lean{Lean4Lean.Pattern.RHS.Uses}
\leanok
\contrib{iota} \srcloc{Lean4Lean/Theory/Typing/Pattern.lean}{170}
\uses{inductive:pattern-rhs}
$r.\mathsf{Uses}\;x$ is defined by recursion on the reduct: false at \texttt{fixed},
equality at \texttt{var}, disjunction at \texttt{app}.
\thesisref{no direct counterpart; support for \texttt{RHS.Generic}}
\end{definition}

\begin{definition}[A generic instance of a rule, \texttt{RHS.Generic}]
\label{def:pattern-rhs-generic}
\lean{Lean4Lean.Pattern.RHS.Generic}
\leanok
\contrib{iota} \srcloc{Lean4Lean/Theory/Typing/Pattern.lean}{181}
\uses{def:pattern-rhs-uses, inductive:pattern-rhs}
$r.\mathsf{Generic}\;m_2\;n$ says that every hole $r$ uses is sent by $m_2$ to a de Bruijn
variable below $n$, that distinct used holes get distinct variables, and that every variable
below $n$ is hit. The holes $r$ does not use are unconstrained.
\thesisref{axioms.tex \S2.6.4 (\texttt{sec:iota}), ``the reduction rule is all substitution
instances of this rule for all the variables left of the turnstile''}
\reviewnote{This is the modelling idea behind \texttt{VEnv.PatTyped}: the thesis's schematic
rule, stated in the context $\Gamma, C{:}\kappa, e{::}\varepsilon, b{::}\beta$, becomes one
typing obligation over a context of exactly the used holes. The docstring is candid that an
$\iota$ rule's index arguments and constructor parameters are tied to nothing by it.}
\end{definition}

\begin{definition}[A positive-occurrence form of the side conditions, \texttt{Check.Realizes}]
\label{def:pattern-check-realizes}
\lean{Lean4Lean.Pattern.Check.Realizes}
\leanok
\contrib{iota} \srcloc{Lean4Lean/Theory/Typing/Pattern.lean}{307}
\uses{inductive:pattern-rhs}
$ck.\mathsf{Realizes}\;m_1\;m_2\;chk$ says the list $chk$ of triples enumerates exactly the
\texttt{defeq} entries of $ck$, with each triple's first two components equal to the
corresponding instantiated sides; the third component is the type at which the two sides are
to be related.
\thesisref{no direct counterpart; the device that makes the \S2.6.4 rule an inductive
constructor}
\reviewnote{A good design move: \texttt{Check.OK} mentions the definitional-equality relation
and so cannot occur positively in the \texttt{IsDefEq} inductive, whereas \texttt{Realizes}
is a pure predicate, which is what makes the new \texttt{IsDefEq.pat} constructor legal.}
\end{definition}

\begin{theorem}[Realizer to \texttt{Check.OK}]
\label{thm:pattern-realizes-took}
\lean{Lean4Lean.Pattern.Check.Realizes.toOK}
\leanok
\contrib{iota} \srcloc{Lean4Lean/Theory/Typing/Pattern.lean}{318}
\uses{def:pattern-check-realizes, inductive:pattern-rhs}
If $chk$ realizes $ck$ at $(m_1, m_2)$ and every triple of $chk$ has its first two components
related by \texttt{defeq}, then $ck.\mathsf{OK}\;\mathsf{defeq}\;m_1\;m_2$.
\thesisref{no direct counterpart; interface to the abstract development}
\end{theorem}
\begin{proof}
\leanok
Induction on $ck$, destructing $chk$ in step and rewriting each side by the realizer's two
equations. It is what feeds the premises of the \texttt{pat} rule into the
\texttt{Check.OK}-based Church-Rosser development.
\end{proof}

\begin{theorem}[\texttt{Check.OK} to realizer]
\label{thm:pattern-ok-exists-realizer}
\lean{Lean4Lean.Pattern.Check.OK.exists_realizer}
\leanok
\contrib{iota} \srcloc{Lean4Lean/Theory/Typing/Pattern.lean}{332}
\uses{def:pattern-check-realizes}
From $ck.\mathsf{OK}\;(\lambda a\,b.\;\exists t,\;rel\;a\;b\;t)\;m_1\;m_2$ one obtains a
realizer $chk$ for $ck$ together with $rel$ at each of its triples.
\thesisref{no direct counterpart; interface to the abstract development}
\end{theorem}
\begin{proof}
\leanok
Induction on $ck$, choosing the type witness out of each existential and consing the triple
onto the realizer. Together with \ref{thm:pattern-realizes-took} this is exactly the
interface \texttt{VEnv.toParams} needs and nothing more.
\end{proof}

\begin{lemma}[Transport lemmas for the pat reduction rule]
\label{family:pattern-transport-iota}
\lean{Lean4Lean.Pattern.RHS.apply_closedN, Lean4Lean.Pattern.Matches.closedN, Lean4Lean.Pattern.RHS.apply_levelWF, Lean4Lean.Pattern.Matches.levelWF, Lean4Lean.Pattern.RHS.instL_apply, Lean4Lean.Pattern.matches_instL, Lean4Lean.Pattern.Check.Realizes.map_liftN, Lean4Lean.Pattern.Check.Realizes.map_instN, Lean4Lean.Pattern.Check.Realizes.map_instL, Lean4Lean.Pattern.RHS.subst_apply, Lean4Lean.Pattern.matches_subst, Lean4Lean.Pattern.Check.Realizes.map_subst}
\leanok
\contrib{iota} \srcloc{Lean4Lean/Theory/Typing/Pattern.lean}{354}
\uses{inductive:pattern-rhs, inductive:pattern-matches, def:pattern-check-realizes, def:closedn, def:instl, fam:subst-basic}
A reduct is \texttt{ClosedN k} as soon as every hole is, and \texttt{LevelWF U} as soon as
every hole is and every match level is \texttt{WF U}; conversely matching a closed,
respectively level-well-formed, expression yields closed, respectively level-well-formed,
holes (and \texttt{WF U} levels). \texttt{RHS.apply} and \texttt{Matches} commute with level
instantiation and with arbitrary substitution; \texttt{Realizes} commutes with all four of
lifting, instantiation, level instantiation and substitution (its lifting and instantiation
counterparts for \texttt{apply}/\texttt{Matches} being \ref{family:pattern-transport-master}).
\thesisref{no direct counterpart; regularity plumbing for the \S2.6.4 rule}
\end{lemma}
\begin{proof}
\leanok
Each is a structural induction on the reduct, on the \texttt{Matches} derivation or on the
\texttt{Check}, with function extensionality on the hole assignments. The twelve lemmas are
what make the \texttt{pat} case of every \texttt{IsDefEq} recursion mechanical, and they are
consumed by \texttt{Theory/Typing/Lemmas.lean} and \texttt{Strong.lean}.
\end{proof}

\section{The iota reduct}

\begin{definition}[The path of the i-th argument of a spine, \texttt{varN\_pathOf}]
\label{def:pattern-varn-pathof}
\lean{Lean4Lean.Pattern.varN_pathOf}
\leanok
\contrib{iota} \srcloc{Lean4Lean/Theory/Typing/Pattern.lean}{494}
\uses{inductive:pattern-core, inductive:pattern-matches}
$\texttt{varN\_pathOf}\;k\;i\;h$ selects argument $i$ of a $k$-ary spine as the path
$\mathsf{some}^{k-1-i}\,\mathsf{none}$. It is defined by recursion on $k$, the zero case
being discharged by the absurd hypothesis $i < 0$.
\thesisref{no direct counterpart; addressing of the \S2.6.4 redex arguments}
\end{definition}

\begin{definition}[The holes an iota reduct uses, \texttt{iotaPaths}]
\label{def:pattern-iota-paths}
\lean{Lean4Lean.SimplePattern.iotaPaths}
\leanok
\contrib{iota} \srcloc{Lean4Lean/Theory/Typing/Pattern.lean}{535}
\uses{def:pattern-varn-pathof, inductive:simple-pattern}
$\texttt{iotaPaths}\;r\;c\;k\;nind\;cnp\;nf$ lists the recursor spine's first $k$ argument
holes, of $k+nind$, followed by the constructor spine's last $nf$ holes, of $cnp+nf$. The
slices are the kernel's: $[0, \texttt{getFirstIndexIdx})$ of the recursor arguments and the
trailing fields of the constructor.
\thesisref{axioms.tex \S2.6.4 (\texttt{sec:iota}), the arguments $e$ and $b$ of $e_c\;b\;v$}
\reviewnote{The two \texttt{pmap}s with their \texttt{omega} side proofs are heavy for what
they say, but correct.}
\end{definition}

\begin{definition}[The iota-rule reduct as a pattern RHS, \texttt{iotaRHS}]
\label{def:pattern-iota-rhs}
\lean{Lean4Lean.SimplePattern.iotaRHS', Lean4Lean.SimplePattern.iotaRHS}
\leanok
\contrib{iota} \srcloc{Lean4Lean/Theory/Typing/Pattern.lean}{569}
\uses{def:pattern-iota-paths, inductive:pattern-rhs}
$\texttt{iotaRHS'}\;r\;c\;k\;nind\;cnp\;nf\;rhs\;h$ is the closed template $rhs$, as a
\texttt{fixed} reduct, applied left-nested to the holes of \texttt{iotaPaths};
\texttt{iotaRHS} specialises $k$ to $np+nm+nmin$ for a recursor with $np$ parameters, $nm$
motives, $nmin$ minors and $nind$ indices firing on a constructor with $cnp$ parameters and
$nf$ fields.
\thesisref{axioms.tex \S2.6.4 (\texttt{sec:iota}), the right-hand side $e_c\;b\;v$, presented
as Lean's stored rule template applied to the recursor prefix and the fields}
\reviewnote{Splitting the shape reduction sees (\texttt{iotaRHS'}) from the typed split
(\texttt{iotaRHS}) is a good separation, and the $cnp$ versus $np$ distinction anticipates
the auxiliary recursors of nested inductives even though \texttt{VInductDecl.WF} currently
excludes them.}
\end{definition}

\begin{lemma}[Instantiating a template applied to holes]
\label{thm:pattern-apply-foldl-var}
\lean{Lean4Lean.Pattern.RHS.apply_foldl_var}
\leanok
\contrib{trproj} \srcloc{Lean4Lean/Theory/Typing/Pattern.lean}{588}
\uses{inductive:pattern-rhs, def:mkapps}
Instantiating $(l.\mathsf{map}\;\mathsf{RHS.var}).\mathsf{foldl}\;\mathsf{RHS.app}\;f$ at a
match gives $(f.\mathsf{apply}\;m_1\;m_2).\mathsf{mkApps}\,(l.\mathsf{map}\;m_2)$.
\thesisref{no direct counterpart; computation lemma for \S2.6.4 reducts}
\end{lemma}
\begin{proof}
\leanok
Induction on the hole list, unfolding \texttt{foldl} and \texttt{mkApps} one step at a time.
\reviewnote{The line blame tags this \texttt{trproj}, but the code is byte-identical on
\texttt{iota} (ported as commit \texttt{3de1dcb}); it belongs to the $\iota$ story, not the
projection story.}
\end{proof}

\begin{theorem}[A saturated constant application matches a spine pattern]
\label{thm:pattern-matches-varn-const}
\lean{Lean4Lean.Pattern.matches_varN_const}
\leanok
\contrib{trproj} \srcloc{Lean4Lean/Theory/Typing/Pattern.lean}{598}
\uses{def:pattern-varn-pathof, inductive:pattern-matches, def:mkapps}
For every $n$ and every argument list of length $n$, the expression
$(\mathsf{const}\;c\;ls).\mathsf{mkApps}\;args$ matches
$(\mathsf{const}\;c).\mathsf{varN}\;n$ at levels $ls$, with a hole assignment sending
$\texttt{varN\_pathOf}\;n\;i$ to $args[i]$.
\thesisref{axioms.tex \S2.6.4 (\texttt{sec:iota}), recognising a redex}
\end{theorem}
\begin{proof}
\leanok
Induction on $n$, peeling the last argument off with \texttt{List.eq\_nil\_or\_concat} and
computing \texttt{varN\_pathOf} at the head by \texttt{dif\_pos}/\texttt{dif\_neg}. It is
what lets the verified WHNF (\srcloc{Lean4Lean/Verify/TypeChecker/WHNF.lean}{108}) and
\srcloc{Lean4Lean/Verify/Environment/Lemmas.lean}{1099} present a kernel $\iota$ redex as a
pattern match. The line blame tags this \texttt{trproj}; the code is identical on
\texttt{iota}.
\end{proof}

\begin{theorem}[The iota reduct on a concrete match]
\label{thm:pattern-iota-rhs-apply}
\lean{Lean4Lean.SimplePattern.iotaRHS'_apply}
\leanok
\contrib{trproj} \srcloc{Lean4Lean/Theory/Typing/Pattern.lean}{625}
\uses{def:pattern-iota-paths, def:pattern-iota-rhs, def:mkapps, def:instl}
If the recursor spine's holes are $recArgs$ of length $k+nind$ and the constructor spine's
are $ctorArgs$ of length $cnp+nf$, then
instantiating the reduct gives $(rhs.\mathsf{instL}\;m_1).\mathsf{mkApps}\,(recArgs.\mathsf{take}\;k \mathbin{+\!\!+} ctorArgs.\mathsf{drop}\;cnp)$.
\thesisref{axioms.tex \S2.6.4 (\texttt{sec:iota}), the reduct of a closed instance of the
rule}
\end{theorem}
\begin{proof}
\uses{thm:pattern-apply-foldl-var}
\leanok
Unfold \texttt{iotaRHS'}, apply \ref{thm:pattern-apply-foldl-var}, rewrite the \texttt{map}
over the two \texttt{pmap}ed ranges, and compare element-wise with
\texttt{List.ext\_getElem} against \texttt{take} and \texttt{drop}. This is the theorem that
makes the model's $\iota$ reduct literally equal to the kernel's
\texttt{inductiveReduceRecCore} reduct, and it is the pivot of that function's correctness
proof. The line blame tags it \texttt{trproj}; the code is identical on \texttt{iota}.
\end{proof}

\begin{lemma}[What an iota reduct retains of the telescope]
\label{family:pattern-iota-counts}
\lean{Lean4Lean.Pattern.RHS.spine, Lean4Lean.Pattern.RHS.spine_foldl_var, Lean4Lean.Pattern.RHS.iotaCounts, Lean4Lean.SimplePattern.iotaPaths_countP_isLeft, Lean4Lean.SimplePattern.iotaPaths_countP_isRight, Lean4Lean.SimplePattern.iotaRHS'_spine, Lean4Lean.SimplePattern.iotaRHS'_iotaCounts, Lean4Lean.SimplePattern.iotaRHS_iotaCounts}
\leanok
\contrib{iota} \srcloc{Lean4Lean/Theory/Typing/Pattern.lean}{511}
\uses{def:pattern-iota-rhs, def:pattern-iota-paths}
A reduct of the form \texttt{fixed c} applied left-nested to \texttt{var} holes decomposes
into head and hole list (\texttt{spine}); for an $\iota$ reduct this yields the template
together with the numbers of recursor-side and constructor-side holes
(\texttt{iotaCounts}), which for \texttt{iotaRHS} are the template, $np+nm+nmin$ and $nf$.
So the reduct depends on the motive/minor split only through their sum.
\thesisref{axioms.tex \S2.6.4 (\texttt{sec:iota}), commentary on what $e_c\;b\;v$ retains}
\end{lemma}
\begin{proof}
\leanok
\texttt{spine\_foldl\_var} is an induction on the hole list; the two counting lemmas are
\texttt{List.countP} computations over the \texttt{pmap}ed ranges of \texttt{iotaPaths}.
\reviewnote{None of these eight declarations is referenced outside
\texttt{Pattern.lean}, and inside it they are used only by each other: about seventy lines
of purely documentary material, better expressed as a doc comment plus one \texttt{example}.}
\end{proof}

\begin{definition}[Template-headed reducts, \texttt{TemplateHeaded}]
\label{inductive:pattern-template-headed}
\lean{Lean4Lean.Pattern.RHS.TemplateHeaded, Lean4Lean.Pattern.RHS.TemplateHeaded.foldl_app, Lean4Lean.Pattern.RHS.TemplateHeaded.apply_lam_or_app, Lean4Lean.Pattern.RHS.TemplateHeaded.apply_ne_forallE, Lean4Lean.Pattern.RHS.TemplateHeaded.apply_ne_sort, Lean4Lean.SimplePattern.iotaRHS'_templateHeaded, Lean4Lean.SimplePattern.iotaRHS_templateHeaded}
\leanok
\contrib{iota} \srcloc{Lean4Lean/Theory/Typing/Pattern.lean}{690}
\uses{inductive:pattern-rhs, def:pattern-iota-rhs}
A reduct is template-headed when it is a \texttt{fixed} closed $\lambda$-abstraction applied
to any arguments. Every instance of such a reduct is an application or an abstraction
(a two-case induction on the derivation), hence never a sort and never a $\Pi$-type; and an
$\iota$ reduct over a $\lambda$-template is template-headed. The predicate is defined and
proved on its own; it is \texttt{PatWF} (\ref{def:patwf}) that is built from it, not the
other way round.
\thesisref{no direct counterpart; the ``a reduction rule computes'' half of
\texttt{VEnv.PatWF}}
\reviewnote{This is what separates a registered reduction rule from a definitional axiom: it
is used in \texttt{Theory/Typing/Lemmas.lean} to show a registered rule cannot rewrite a sort
or a $\Pi$-type, which the inversion lemmas of the type system need. A well-judged
invariant.}
\end{definition}

\section{Registering an inductive block}

\begin{definition}[Registering one iota rule, \texttt{addRecRule}]
\label{def:ind-add-rec-rule}
\lean{Lean4Lean.VEnv.addRecRule}
\leanok
\contrib{iota} \srcloc{Lean4Lean/Theory/Inductive.lean}{257}
\uses{def:pattern-iota-rhs, def:venv-addpat, inductive:simple-pattern, def:getmajoridx, struct:vrecursor, struct:vrecrule}
$env.\mathsf{addRecRule}\;r\;ru$ registers, when $ru.rhs$ is closed, the pattern
$\mathsf{iota}\;r.name\;M\;ru.ctor\;N$ with $M$ the major index and
$N = ru.\mathsf{ctorParams} + ru.\mathsf{nfields}$, with reduct \texttt{iotaRHS} and the
trivial side condition \texttt{Check.true}; otherwise it fails.
\thesisref{axioms.tex \S2.6.4 (\texttt{sec:iota}), the rule
$\mathrm{rec}_P\;C\;e\;p[b]\;(c\;b) \rightsquigarrow e_c\;b\;v$}
\reviewnote{Only the constructor rule of \S2.6.4 is registered. K-like reduction, the
section's second rule, which the kernel does perform (\texttt{toCtorWhenK},
\srcloc{Lean4Lean/Inductive/Reduce.lean}{29}, called from
\srcloc{Lean4Lean/Inductive/Reduce.lean}{110}), is not modelled, and \texttt{VRecursor.k}
is recorded but unused. Style nit: the major index is spelled out here as
\texttt{r.numParams + r.numMotives + r.numMinors + r.numIndices} while every lemma about it
uses \texttt{r.getMajorIdx}.}
\end{definition}

\begin{definition}[The four stages of \texttt{addInduct}]
\label{family:ind-add-stages}
\lean{Lean4Lean.VInductDecl.addTypes, Lean4Lean.VInductDecl.addCtors, Lean4Lean.VInductDecl.addRecs, Lean4Lean.VInductDecl.addRules, Lean4Lean.VInductDecl.addTypesCtors, Lean4Lean.VInductDecl.addTypesCtorsRecs, Lean4Lean.VInductDecl.consts}
\leanok
\contrib{iota} \srcloc{Lean4Lean/Theory/Inductive.lean}{280}
\uses{def:ind-add-rec-rule, def:venv-basic, struct:vinductdecl}
Four \texttt{Option}-monadic folds add, in turn, the type formers as constants, all
constructors in block order, the recursors, and every recursor rule as an $\iota$ rule; the
composites \texttt{addTypesCtors} and \texttt{addTypesCtorsRecs} name the prefixes, and
\texttt{consts} lists the block's name-constant pairs in stage order.
\thesisref{axioms.tex \S2.6.1 (\texttt{sec:inductive}), ``$\mu t{:}F.\,K$ and its
constructors are implemented as additional axiomatic constant symbols''}
\reviewnote{The staging exists precisely so that \texttt{VInductDecl.WF} can type each kind
of constant in the environment the kernel checks it in, and the docstring is honest that the
kernel interleaves recursors with their rules, so only the resulting environment agrees, not
the order.}
\end{definition}

\begin{definition}[Extending an environment with an inductive block, \texttt{addInduct}]
\label{def:ind-add-induct}
\lean{Lean4Lean.VEnv.addInduct}
\leanok
\contrib{iota} \srcloc{Lean4Lean/Theory/Inductive.lean}{316}
\uses{family:ind-add-stages, struct:venv}
$env.\mathsf{addInduct}\;decl$ is the chain
$\mathsf{addTypes} \gg\!\!= \mathsf{addCtors} \gg\!\!= \mathsf{addRecs} \gg\!\!= \mathsf{addRules}$,
returning \texttt{none} on a name clash or a non-closed rule reduct.
\thesisref{axioms.tex \S2.6.1 (\texttt{sec:inductive}) through \S2.6.4 (\texttt{sec:iota})}
\reviewnote{On \texttt{master} this was \texttt{def VEnv.addInduct := sorry}, an opaque
\texttt{Option VEnv}, so no downstream statement mentioning it could be proved or refuted.}
\end{definition}

\begin{definition}[Well-formedness of a direct mutual inductive block, \texttt{VInductDecl.WF}]
\label{structure:ind-wf}
\lean{Lean4Lean.VInductDecl.WF}
\leanok
\contrib{mixed} \srcloc{Lean4Lean/Theory/Inductive.lean}{335}
\uses{def:ind-rec-shape, def:ind-rule-shape, def:ind-ctor-result, def:ind-ctor-positive, def:ind-large-elim, def:ind-major-former, def:ind-minor-for, def:ind-field-ctx, family:ind-add-stages, def:pattyped, def:pattern-iota-rhs, inductive:simple-pattern, def:getmajoridx, def:hastype-istype, fam:vlevel-lattice}
A record of twenty clauses. The type formers, constructors and recursors are typed in the
stage in which the kernel checks them (\texttt{types\_wf}, \texttt{ctors\_wf},
\texttt{recs\_wf}); universe parameters are shared; \texttt{universes} asks for one sort
$\ell$ that is the result sort of every type former of the block and dominates every
constructor field, $\mathrm{imax}(u,\ell) \le \ell$, with \texttt{LargeElim} whenever some
recursor takes the extra universe; \texttt{recs\_elim} encodes $\kappa$ by pinning the motive
sorts; \texttt{rec\_params} and \texttt{rec\_counts} fix the telescope split (one motive per
type former, one minor per constructor of the block, indices as many as the eliminated
former has) and \texttt{rec\_shape} the telescope itself; constructors are
\texttt{CtorResult} and \texttt{CtorPositive} with parameters shared with their former; the
block is direct (\texttt{recs\_over\_block}, \texttt{rules\_ctor},
\texttt{types\_have\_rec}); each recursor has exactly one rule per constructor of the former
it eliminates (\texttt{rules\_nodup}, \texttt{rules\_total}); each reduct has
\texttt{RuleShape} against the matching minor (\texttt{rule\_shape}); and each $\iota$ rule
is \texttt{PatTyped} in the stage-2 environment (\texttt{rules\_wf}).
The attribution is mixed only in a trivial sense: the line blame tags the
\texttt{types\_have\_rec} field \texttt{trproj}, but that is a line re-wrap and the field
exists identically on \texttt{iota}.
\thesisref{axioms.tex \S2.6.1 (\texttt{sec:inductive}), \S2.6.2 (\texttt{sec:large\_elim}),
\S2.6.3 and \S2.6.4 (\texttt{sec:iota})}
\reviewnote{This replaces \texttt{master}'s \texttt{def VInductDecl.WF := sorry} and is the
centrepiece of the contribution. Four limitations are stated in the structure's own
docstring: nested inductives are excluded; K-like reduction and structure $\eta$ are not
modelled; and positivity and \texttt{universes} read manifest binders where the kernel
\texttt{whnf}s. \texttt{rule\_shape}'s field docstring adds a fifth, that the recursive
arguments are pinned in number only. Two further ones are documented nowhere:
\texttt{universes} bundles the purely syntactic bound
\texttt{decl.nparams} $\le$ \texttt{t.type.piArity} into a clause guarded by stage 0
succeeding, so that arity fact is unavailable if stage 0 fails; and no clause mentions
\texttt{VRecursor.all}, which is pinned to the kernel's only downstream, by
\texttt{TrRecursor.all} (\srcloc{Lean4Lean/Verify/Environment/Basic.lean}{240}). Above all,
and also undocumented, no theorem derives this record from the kernel's own checker.}
\end{definition}

\begin{lemma}[Consequences of block well-formedness]
\label{family:ind-wf-corollaries}
\lean{Lean4Lean.VInductDecl.WF.ctors_have_rules, Lean4Lean.VInductDecl.WF.rules_own_params}
\leanok
\contrib{iota} \srcloc{Lean4Lean/Theory/Inductive.lean}{430}
\uses{structure:ind-wf}
Every constructor of the block has a rule in some recursor of the block, and every rule
records its recursor's parameter count, $ru.\mathsf{ctorParams} = r.\mathsf{numParams}$.
\thesisref{no direct counterpart; a small interface for the front end}
\end{lemma}
\begin{proof}
\leanok
\texttt{ctors\_have\_rules} chains \texttt{types\_have\_rec} with \texttt{rules\_total};
\texttt{rules\_own\_params} combines \texttt{rules\_ctor}, which says the rule's constructor
is one of the block's and so carries \texttt{decl.nparams}, with \texttt{rec\_params}. Both
are consumed by \texttt{Verify/Environment/Basic.lean}, and
\texttt{rules\_own\_params} is what collapses the $cnp$ versus $np$ distinction of
\texttt{iotaRHS} for a direct block.
\end{proof}

\section{Environment bookkeeping for \texorpdfstring{\texttt{addInduct}}{addInduct}}

This section is \srcloc{Lean4Lean/Theory/Typing/InductiveLemmas.lean}{1}. On \texttt{master}
the file contained only \texttt{theorem addInduct\_WF ... := sorry}.

\begin{lemma}[Generic reasoning principles for monadic folds over environments]
\label{family:indlem-foldlm}
\lean{Lean4Lean.VEnv.foldlM_le, Lean4Lean.VEnv.foldlM_inv, Lean4Lean.VEnv.foldlM_mono_of_mem, Lean4Lean.VEnv.foldlM_pats_preserved, Lean4Lean.VEnv.foldlM_defeqs_preserved, Lean4Lean.VEnv.foldlM_pats_inv_mem, Lean4Lean.VEnv.foldlM_step_success, Lean4Lean.VEnv.foldlM_addConst_ordered}
\leanok
\contrib{iota} \srcloc{Lean4Lean/Theory/Typing/InductiveLemmas.lean}{22}
\uses{struct:venv-le, def:ordered, def:venv-basic}
For a successful fold $l.\mathsf{foldlM}\;f\;init = \mathsf{some}\;r$: monotonicity if each
step is monotone; invariance of a step-preserved predicate; propagation of a property
established at one list element; preservation of \texttt{pats} and of \texttt{defeqs};
membership-tracking inversion of a pattern present after the fold; existence of a successful
step for each element; and orderedness of an \texttt{addConst} fold whose constants are all
well-formed at the start.
\thesisref{no direct counterpart; proof infrastructure}
\end{lemma}
\begin{proof}
\leanok
All by induction on the list, splitting the \texttt{Option.bind} at each step;
\texttt{foldlM\_addConst\_ordered} runs \texttt{foldlM\_inv} with the invariant
``ordered and above the initial environment''.
\reviewnote{These are generic over an arbitrary element type and have nothing to do with
inductives, yet they live in namespace \texttt{Lean4Lean.VEnv} inside a file called
\texttt{InductiveLemmas}.}
\end{proof}

\begin{theorem}[Decomposition of a successful \texttt{addInduct}]
\label{thm:indlem-add-induct-stages}
\lean{Lean4Lean.VEnv.addInduct_stages}
\leanok
\contrib{iota} \srcloc{Lean4Lean/Theory/Typing/InductiveLemmas.lean}{76}
\uses{def:ind-add-induct, family:ind-add-stages}
If $env.\mathsf{addInduct}\;decl = \mathsf{some}\;env'$ then there are intermediate
environments $env_T$, $env_C$, $env_R$ at which the four stages each succeed.
\thesisref{no direct counterpart; proof infrastructure}
\end{theorem}
\begin{proof}
\leanok
Repeated \texttt{Option.bind\_eq\_some\_iff} after unfolding the composite stage
definitions. It is the workhorse of almost every other lemma in the file and of the front
end.
\end{proof}

\begin{lemma}[Adding an inductive only grows the environment]
\label{family:indlem-stage-le}
\lean{Lean4Lean.VEnv.addRecRule_le, Lean4Lean.VEnv.addTypes_le, Lean4Lean.VEnv.addCtors_le, Lean4Lean.VEnv.addRecs_le, Lean4Lean.VEnv.addRules_le, Lean4Lean.VEnv.addInduct_le}
\leanok
\contrib{iota} \srcloc{Lean4Lean/Theory/Typing/InductiveLemmas.lean}{32}
\uses{family:ind-add-stages, struct:venv-le}
Each stage of \texttt{addInduct}, and \texttt{addInduct} itself, yields an environment above
the one it started from.
\thesisref{no direct counterpart; proof infrastructure}
\end{lemma}
\begin{proof}
\uses{family:indlem-foldlm, thm:indlem-add-induct-stages, fam:addpat}
\leanok
\texttt{foldlM\_le} over \texttt{addConst\_le} and \texttt{addPat\_le}, chained through
\ref{thm:indlem-add-induct-stages} for the composite.
\end{proof}

\begin{lemma}[Every recursor rule is registered as an iota rule]
\label{family:indlem-pat-registration}
\lean{Lean4Lean.VEnv.addRecRule_pats, Lean4Lean.VEnv.addRules_pat, Lean4Lean.VEnv.addInduct_pat}
\leanok
\contrib{iota} \srcloc{Lean4Lean/Theory/Typing/InductiveLemmas.lean}{41}
\uses{def:ind-add-rec-rule, def:pattern-iota-rhs, inductive:simple-pattern}
After \texttt{addRecRule}, \texttt{addRules} or \texttt{addInduct}, for every recursor $r$ of
the block and every rule $ru$ of $r$ with closed reduct, the entry keyed by
$\mathsf{iota}\;r.name\;r.\mathsf{getMajorIdx}\;ru.ctor\;(ru.\mathsf{ctorParams}+ru.\mathsf{nfields})$
with reduct \texttt{iotaRHS} and check \texttt{true} is present in \texttt{pats}.
\thesisref{axioms.tex \S2.6.4 (\texttt{sec:iota})}
\end{lemma}
\begin{proof}
\uses{family:indlem-foldlm, thm:indlem-add-induct-stages, fam:addpat}
\leanok
\texttt{addPat\_self} for the single step, then \texttt{foldlM\_mono\_of\_mem} twice, over
the recursors and over their rules, to carry the entry through the remaining fold, and
\ref{thm:indlem-add-induct-stages} to reach it from \texttt{addInduct}. \texttt{addInduct\_pat}
is the interface the verified type checker uses to know that a rule it looked up in the
kernel is really registered in the model.
\end{proof}

\begin{lemma}[Each stage preserves orderedness]
\label{family:indlem-stage-ordered}
\lean{Lean4Lean.VEnv.addTypes_ordered, Lean4Lean.VEnv.addCtors_ordered, Lean4Lean.VEnv.addRecs_ordered, Lean4Lean.VEnv.addTypesCtors_ordered, Lean4Lean.VEnv.addTypesCtorsRecs_ordered, Lean4Lean.VEnv.addRules_ordered}
\leanok
\contrib{iota} \srcloc{Lean4Lean/Theory/Typing/InductiveLemmas.lean}{171}
\uses{structure:ind-wf, def:ordered, family:ind-add-stages}
Stages 0 to 2 preserve \texttt{Ordered} because the corresponding clause of
\texttt{VInductDecl.WF} types the constants in exactly the environment the stage starts
from; stage 3 preserves it because every registration is an \texttt{Ordered.pat} step.
\thesisref{no direct counterpart; the admissibility half of \S2.6.1 to \S2.6.4}
\end{lemma}
\begin{proof}
\uses{family:indlem-foldlm, inductive:pattern-template-headed, family:ind-shape-corollaries, def:patwf, fam:patwf-mono}
\leanok
\texttt{foldlM\_addConst\_ordered} for the three constant stages. For
\texttt{addRules\_ordered}, a doubly nested \texttt{foldlM\_inv} with invariant
``ordered and above $env_R$'': the \texttt{PatWF} obligation is \texttt{rules\_wf} moved
forward by \texttt{PatTyped.mono}, together with \texttt{RuleShape.lam} and
\texttt{iotaRHS\_templateHeaded} for the template shape. This is where the staging of
\texttt{VInductDecl.WF} pays off: the three constant stages are one to three lines each.
\end{proof}

\begin{theorem}[Soundness of \texttt{addInduct}]
\label{thm:indlem-add-induct-wf}
\lean{Lean4Lean.VEnv.addInduct_WF}
\leanok
\contrib{mixed} \srcloc{Lean4Lean/Theory/Typing/InductiveLemmas.lean}{229}
\uses{structure:ind-wf, def:ind-add-induct, def:ordered}
Extending an \texttt{Ordered} environment with an inductive declaration well-formed in it
yields an \texttt{Ordered} environment. The statement line is \texttt{master}'s; the record
it refers to and the proof are the \texttt{iota} branch's.
\thesisref{axioms.tex \S2.6.1 (\texttt{sec:inductive}) through \S2.6.4 (\texttt{sec:iota}),
admissibility of an inductive block}
\end{theorem}
\begin{proof}
\uses{family:indlem-stage-ordered}
\leanok
Split the final \texttt{Option.bind} of \texttt{addInduct}, then apply
\texttt{addTypesCtorsRecs\_ordered} followed by \texttt{addRules\_ordered}.
\reviewnote{This is the headline result of the branch for this file: \texttt{master}
declared exactly this statement with \texttt{sorry}. It is proved only relative to the new
\texttt{VInductDecl.WF}, which is itself part of the contribution and remains an assumption
on the front-end side.}
\end{proof}

\begin{lemma}[Which extensions touch pats and which touch defeqs]
\label{family:indlem-pats-defeqs-preserved}
\lean{Lean4Lean.VEnv.addConst_pats, Lean4Lean.VEnv.addConst_defeqs, Lean4Lean.VEnv.addDefEq_pats, Lean4Lean.VEnv.addConsts_pats, Lean4Lean.VEnv.addConsts_defeqs, Lean4Lean.VEnv.addDefEqs_pats, Lean4Lean.VEnv.addDefEqs_le, Lean4Lean.VEnv.addQuot_pats, Lean4Lean.VEnv.addQuot_le, Lean4Lean.VEnv.addTypes_pats, Lean4Lean.VEnv.addCtors_pats, Lean4Lean.VEnv.addRecs_pats, Lean4Lean.VEnv.addTypes_defeqs, Lean4Lean.VEnv.addCtors_defeqs, Lean4Lean.VEnv.addRecs_defeqs, Lean4Lean.VEnv.addRecRule_defeqs, Lean4Lean.VEnv.addRules_defeqs, Lean4Lean.VEnv.addInduct_defeqs}
\leanok
\contrib{iota} \srcloc{Lean4Lean/Theory/Typing/InductiveLemmas.lean}{246}
\uses{def:venv-basic, def:vdecl-wf, family:ind-add-stages}
\texttt{addConst}, \texttt{addConsts}, \texttt{addDefEq}, \texttt{addDefEqs},
\texttt{addQuot} and the three constant stages of \texttt{addInduct} leave \texttt{pats}
unchanged; \texttt{addConst}, \texttt{addConsts} and all four stages leave \texttt{defeqs}
unchanged. In particular an inductive block contributes $\iota$ rules and never a
definitional axiom.
\thesisref{no direct counterpart; the separation \S2.6.4 needs from \S2.5}
\end{lemma}
\begin{proof}
\uses{family:indlem-foldlm, thm:indlem-add-induct-stages}
\leanok
Case analysis on \texttt{addConst} followed by \texttt{foldlM\_pats\_preserved} or
\texttt{foldlM\_defeqs\_preserved}, with \ref{thm:indlem-add-induct-stages} for the
composites. This is exactly the separation the \texttt{Params} instance needs.
\reviewnote{\texttt{addQuot\_le}, \texttt{addQuot\_pats} and the \texttt{addDefEqs} lemmas
are about other declaration kinds and sit oddly in a file called \texttt{InductiveLemmas};
the three \texttt{defeqs} lemmas for rules are tagged \texttt{trproj} by the blame but exist
identically on \texttt{iota}.}
\end{proof}

\begin{lemma}[Full specification of a successful \texttt{addConst} fold]
\label{family:indlem-addconst-spec}
\lean{Lean4Lean.VEnv.addConst_eq, Lean4Lean.VEnv.addConst_foldlM_fresh, Lean4Lean.VEnv.addConst_comm, Lean4Lean.VEnv.addConst_foldlM_perm, Lean4Lean.VEnv.addConst_foldlM_constants_inv, Lean4Lean.VEnv.addConst_foldlM_find, Lean4Lean.VEnv.addConst_foldlM_inj, Lean4Lean.VEnv.addConst_foldlM_nodup, Lean4Lean.VEnv.nodup_map_inj_on}
\leanok
\contrib{iota} \srcloc{Lean4Lean/Theory/Typing/InductiveLemmas.lean}{325}
\uses{def:venv-basic, struct:venv}
A successful \texttt{addConst} sets exactly its own name, which was fresh. In a successful
fold: every registered name was fresh at the start; every name is bound in the result to its
own constant; the naming function is injective on the list and the names have no duplicates;
any constant bound at the end was bound before or is one of the fold's; two steps commute;
and the result is invariant under permuting the list.
\thesisref{no direct counterpart; proof infrastructure}
\end{lemma}
\begin{proof}
\leanok
Induction on the list throughout. \texttt{addConst\_comm} is a \texttt{VEnv.ext}
computation using that the second name was fresh, and \texttt{addConst\_foldlM\_perm} is an
induction on the permutation using it. The permutation and inversion lemmas are what let the
front end reconcile the kernel's per-type insertion order with the model's four-stage order.
\reviewnote{\texttt{nodup\_map\_inj\_on} is a pure \texttt{List} lemma declared as
\texttt{Lean4Lean.VEnv.nodup\_map\_inj\_on}.}
\end{proof}

\begin{lemma}[Stages 0 to 2 are one \texttt{addConst} fold]
\label{thm:indlem-addtypesctorsrecs-eq}
\lean{Lean4Lean.VInductDecl.addTypesCtorsRecs_eq}
\leanok
\contrib{iota} \srcloc{Lean4Lean/Theory/Typing/InductiveLemmas.lean}{410}
\uses{family:ind-add-stages}
$decl.\mathsf{addTypesCtorsRecs}\;env$ equals the \texttt{addConst} fold over
$decl.\mathsf{consts}$.
\thesisref{no direct counterpart; proof infrastructure}
\end{lemma}
\begin{proof}
\uses{family:indlem-addconst-spec}
\leanok
\texttt{simp} with \texttt{List.foldlM\_append} and \texttt{List.foldlM\_map}, then
\texttt{rfl}. It is what lets the whole \texttt{addConst\_foldlM} toolkit apply to the block
at once.
\end{proof}

\begin{lemma}[What each stage binds]
\label{family:indlem-find-fresh-nodup}
\lean{Lean4Lean.VEnv.addTypes_find, Lean4Lean.VEnv.addCtors_fresh, Lean4Lean.VEnv.addCtors_find, Lean4Lean.VEnv.addRecs_fresh, Lean4Lean.VEnv.addRecs_find, Lean4Lean.VEnv.addRecs_name_inj, Lean4Lean.VEnv.addTypes_nodup, Lean4Lean.VEnv.addCtors_nodup, Lean4Lean.VEnv.addRecs_nodup}
\leanok
\contrib{iota} \srcloc{Lean4Lean/Theory/Typing/InductiveLemmas.lean}{532}
\uses{family:ind-add-stages, struct:venv}
After its own stage, every type former, constructor and recursor of the block is bound to its
own constant; constructors and recursors were unbound beforehand; names within each kind have
no duplicates, and two recursors of the block with the same name coincide.
\thesisref{no direct counterpart; proof infrastructure}
\end{lemma}
\begin{proof}
\uses{family:indlem-addconst-spec}
\leanok
Instances of the \texttt{addConst\_foldlM} family after unfolding the stage definition, with
\texttt{List.mem\_flatMap} for the constructor stage.
\end{proof}

\begin{lemma}[Origin of a registered pattern]
\label{family:indlem-pats-origin}
\lean{Lean4Lean.VEnv.addRecRule_pats_inv', Lean4Lean.VEnv.addInduct_pats_origin', Lean4Lean.VEnv.addInduct_pats_origin}
\leanok
\contrib{iota} \srcloc{Lean4Lean/Theory/Typing/InductiveLemmas.lean}{577}
\uses{def:ind-add-rec-rule, def:ind-add-induct, def:pattern-iota-rhs, inductive:simple-pattern}
A pattern entry present after \texttt{addRecRule} or \texttt{addInduct} is either one that was
already there, or exactly the $\iota$ entry of one rule of one recursor of the block: key
\texttt{SimplePattern.iota} \emph{and} reduct \texttt{SimplePattern.iotaRHS}, both read off
that rule. The unprimed version records only the key.
\thesisref{axioms.tex \S2.6.4 (\texttt{sec:iota}), the converse reading of the rule set}
\end{lemma}
\begin{proof}
\uses{family:indlem-pat-registration, family:indlem-foldlm, family:indlem-pats-defeqs-preserved}
\leanok
Inversion of \texttt{addPat} for the single step; then \texttt{foldlM\_pats\_inv\_mem} nested
twice with the per-recursor and per-rule motives spelled out, and the three constant stages
rewritten away by the \texttt{pats}-preservation lemmas.
\reviewnote{The dependent-equality formulation is heavy, but it is what pins the reduct and
not only the key; the primed and unprimed pair is a sensible interface split.}
\end{proof}

\begin{lemma}[The recursors of a block, seen from \texttt{addInduct}]
\label{family:indlem-addinduct-consts}
\lean{Lean4Lean.VEnv.addInduct_rec_fresh, Lean4Lean.VEnv.addInduct_rec_find, Lean4Lean.VEnv.addInduct_recs_name_inj}
\leanok
\contrib{iota} \srcloc{Lean4Lean/Theory/Typing/InductiveLemmas.lean}{660}
\uses{def:ind-add-induct, struct:vrecursor}
A recursor of \texttt{decl} is not already bound in \texttt{env}, is bound to its own constant
in the result, and two recursors of the block sharing a name are equal.
\thesisref{no direct counterpart; proof infrastructure}
\end{lemma}
\begin{proof}
\uses{thm:indlem-add-induct-stages, family:indlem-find-fresh-nodup, family:indlem-stage-le}
\leanok
Thread \ref{thm:indlem-add-induct-stages} and the stage lemmas through monotonicity of the
environment order. These are exactly the facts \ref{thm:indparams-patsiota-induct} needs to
rule out a name clash between an old registered recursor and a new one.
\end{proof}

\begin{theorem}[The constructor an iota rule fires on is registered with the right arity]
\label{thm:indlem-rules-ctor-shape}
\lean{Lean4Lean.VInductDecl.WF.rules_ctor_shape, Lean4Lean.VEnv.addInduct_rule_ctor}
\leanok
\contrib{iota} \srcloc{Lean4Lean/Theory/Typing/InductiveLemmas.lean}{687}
\uses{structure:ind-wf, def:ind-ctor-shape}
For a well-formed block, the constructor of every recursor rule is registered in the stage-1
environment, and hence in the result of \texttt{addInduct}, with a type of
$\mathsf{CtorShape}\,(ru.\mathsf{ctorParams} + ru.\mathsf{nfields})$.
\thesisref{no direct counterpart; the environment-level residue of \S2.6.1}
\end{theorem}
\begin{proof}
\uses{family:ind-shape-corollaries, family:indlem-find-fresh-nodup, thm:indlem-add-induct-stages}
\leanok
Read \texttt{rules\_ctor} to find the constructor in the block, look it up with
\texttt{addCtors\_find}, and convert \texttt{CtorResult} to \texttt{CtorShape} with
\texttt{CtorResult.ctorShape} after rewriting $ru.\mathsf{ctorParams} = decl.\mathsf{nparams}$;
the second statement carries the result forward along the remaining stages. This is the
bridge that ultimately separates recursor heads from constructor heads.
\end{proof}

\begin{theorem}[Injectivity of the iota redex shape]
\label{thm:indlem-iota-topattern-inj}
\lean{Lean4Lean.VEnv.iota_toPattern_inj}
\leanok
\contrib{iota} \srcloc{Lean4Lean/Theory/Typing/InductiveLemmas.lean}{713}
\uses{inductive:simple-pattern}
If two $\mathsf{iota}$ redexes have the same underlying \texttt{Pattern}, their recursor
names, spine arities, constructor names and constructor arities agree.
\thesisref{no direct counterpart; proof infrastructure}
\end{theorem}
\begin{proof}
\uses{family:pattern-spine-combinatorics}
\leanok
Injection on the \texttt{app} node, then \texttt{Pattern.varN\_const\_inj} on each spine.
Small, but used throughout \texttt{InductiveParams} and in the front end.
\end{proof}

\begin{lemma}[Registered reducts are closed]
\label{family:indlem-closed}
\lean{Lean4Lean.VEnv.addRecRule_closed, Lean4Lean.VEnv.addRules_closed}
\leanok
\contrib{iota} \srcloc{Lean4Lean/Theory/Typing/InductiveLemmas.lean}{737}
\uses{def:ind-add-rec-rule, def:closedn}
A successful \texttt{addRecRule} registered a closed reduct, and after a successful
\texttt{addRules} every rule reduct of the declaration is closed.
\thesisref{no direct counterpart; proof infrastructure}
\end{lemma}
\begin{proof}
\uses{family:indlem-foldlm}
\leanok
Inversion of the \texttt{if} in \texttt{addRecRule}, lifted through
\texttt{foldlM\_step\_success} twice. It is what lets the front end produce the closedness
argument of \texttt{iotaRHS} without \texttt{VInductDecl.WF} carrying it.
\end{proof}

\section{A Params instance from the registered rules}

This section is \srcloc{Lean4Lean/Theory/Typing/InductiveParams.lean}{1}. Its goal is to
instantiate the abstract parameter record the Church-Rosser development is stated over with
an environment's own registered $\iota$ rules.

\begin{lemma}[The only application node of an iota redex]
\label{family:indparams-iota-subpattern}
\lean{Lean4Lean.VEnv.app_subpattern_iota', Lean4Lean.VEnv.app_subpattern_iota}
\leanok
\contrib{iota} \srcloc{Lean4Lean/Theory/Typing/InductiveParams.lean}{26}
\uses{inductive:simple-pattern, inductive:pattern-core}
Any application subpattern of $(\mathsf{iota}\;r\;M\;c\;N).\mathsf{toPattern}$ is the
top-level one, whose factors are the recursor spine and the constructor spine.
\thesisref{no direct counterpart; coherence of the rule set}
\end{lemma}
\begin{proof}
\uses{family:pattern-spine-combinatorics}
\leanok
Case analysis on the \texttt{Subpattern} derivation, the \texttt{appL} and \texttt{appR}
cases being excluded by \texttt{not\_app\_subpattern\_varN\_const}.
\end{proof}

\begin{lemma}[Inversion lemmas for pattern intersection at an application]
\label{family:indparams-inter}
\lean{Lean4Lean.Pattern.inter_app_const, Lean4Lean.Pattern.inter_app_var, Lean4Lean.Pattern.inter_app_app}
\leanok
\contrib{iota} \srcloc{Lean4Lean/Theory/Typing/InductiveParams.lean}{43}
\uses{inductive:pattern-core}
An application pattern does not intersect a constant; it intersects a \texttt{var} only
through its function part; and two applications intersect componentwise.
\thesisref{no direct counterpart; coherence of the rule set}
\end{lemma}
\begin{proof}
\leanok
\texttt{simp [Pattern.inter]} after case-splitting on the recursive intersections.
\reviewnote{Three \texttt{Pattern} lemmas declared with \texttt{\_root\_} inside
\texttt{namespace VEnv}; they belong next to \texttt{Pattern.inter} in
\texttt{Pattern.lean}.}
\end{proof}

\begin{lemma}[Only inductive declarations change pats]
\label{family:indparams-pats-eq-or-induct}
\lean{Lean4Lean.VDecl.WF.pats_eq_or_induct', Lean4Lean.VDecl.WF.pats_eq_or_induct}
\leanok
\contrib{iota} \srcloc{Lean4Lean/Theory/Typing/InductiveParams.lean}{69}
\uses{def:vdecl-wf, def:ind-add-induct, structure:ind-wf}
Each \texttt{VDecl.WF} step either leaves \texttt{pats} unchanged (axiom, def, mutual def,
opaque, example, quot) or is the \texttt{addInduct} of a declaration well-formed at that
point.
\thesisref{no direct counterpart; the invariant that \texttt{pats} holds exactly the $\iota$
rules}
\end{lemma}
\begin{proof}
\uses{family:indlem-pats-defeqs-preserved}
\leanok
Case analysis on the constructors of \texttt{VDecl.WF}, each non-inductive case discharged by
the corresponding \texttt{pats}-preservation lemma. It is the induction step of both
\ref{thm:indparams-wf-pats-origin} and \ref{thm:indparams-wf-patsiota}.
\end{proof}

\begin{theorem}[Origin of a pattern along a declaration chain]
\label{thm:indparams-wf-pats-origin}
\lean{Lean4Lean.VEnv.WF'.pats_origin}
\leanok
\contrib{iota} \srcloc{Lean4Lean/Theory/Typing/InductiveParams.lean}{93}
\uses{def:venv-wf, def:ind-add-induct, structure:ind-wf, def:pattern-iota-rhs}
If $env.\mathsf{WF'}\;ds$ and $env.\mathsf{pats}\;p\;rr$, then some suffix of $ds$ begins with
an inductive declaration $decl$ well-formed in the environment $env_0$ reached there, with
$env_0.\mathsf{addInduct}\;decl = \mathsf{some}\;env_1$ and $env_1 \le env$, and $(p, rr)$ is
exactly the $\iota$ entry of one rule of one recursor of $decl$.
\thesisref{axioms.tex \S2.6.4 (\texttt{sec:iota}), locating the declaration a rule came from}
\end{theorem}
\begin{proof}
\uses{family:indparams-pats-eq-or-induct, family:indlem-pats-origin}
\leanok
Induction on the \texttt{WF'} chain, using \texttt{pats\_eq\_or\_induct'} at each step and
\texttt{addInduct\_pats\_origin'} in the inductive case, carrying the environment-order
witness forward.
\reviewnote{Its own docstring calls it the first step of the deferred $\iota$
subject-reduction proof (\texttt{VEnv.WF.patsStrong},
\srcloc{Lean4Lean/Theory/Typing/EnvLemmas.lean}{334}, still \texttt{sorry}). Nothing uses it
yet, so its value is entirely as groundwork for that open obligation.}
\end{proof}

\begin{definition}[The pattern-registry invariant, \texttt{PatsIota}]
\label{structure:indparams-patsiota}
\lean{Lean4Lean.VEnv.PatsIota}
\leanok
\contrib{iota} \srcloc{Lean4Lean/Theory/Typing/InductiveParams.lean}{133}
\uses{inductive:simple-pattern, def:recheaded, def:ctorheaded, struct:venv}
Four fields. \texttt{shape}: every registered pattern is an $\mathsf{iota}$ redex whose
recursor head is a registered constant with a \texttt{RecHeaded} type. \texttt{arity}: the
recursor name determines the spine arity. \texttt{ctor\_shape}: the constructor is a
registered constant with a \texttt{CtorHeaded} type. \texttt{functional}: a pattern
determines its reduct.
\thesisref{no direct counterpart; the population invariant behind the \texttt{Params}
instance}
\reviewnote{Well chosen: exactly strong enough to discharge the five structural
\texttt{Params} conditions and weak enough to be preserved by \texttt{addInduct}. The
\texttt{arity} field is the non-obvious one.}
\end{definition}

\begin{theorem}[A recursor head is never a constructor head]
\label{thm:indparams-patsiota-rec-ne-ctor}
\lean{Lean4Lean.VEnv.PatsIota.rec_ne_ctor}
\leanok
\contrib{iota} \srcloc{Lean4Lean/Theory/Typing/InductiveParams.lean}{148}
\uses{structure:indparams-patsiota, def:recheaded, def:ctorheaded}
Under \texttt{PatsIota}, the recursor name of one registered $\iota$ redex differs from the
constructor name of any registered $\iota$ redex.
\thesisref{no direct counterpart; coherence of the rule set}
\end{theorem}
\begin{proof}
\uses{thm:indlem-iota-topattern-inj}
\leanok
Both names would resolve to the same registered constant, whose type would be both
\texttt{RecHeaded} and \texttt{CtorHeaded}, contradicting \texttt{RecHeaded.not\_ctorHeaded}.
\reviewnote{The separation rests on the type shapes of the constants, not on any tagging in
\texttt{VEnv.constants}, which is precisely why the same trick does not extend to $\delta$
rules.}
\end{proof}

\begin{lemma}[\texttt{PatsIota} survives a non-inductive step]
\label{thm:indparams-patsiota-of-le}
\lean{Lean4Lean.VEnv.PatsIota.of_le}
\leanok
\contrib{iota} \srcloc{Lean4Lean/Theory/Typing/InductiveParams.lean}{161}
\uses{structure:indparams-patsiota, struct:venv-le}
\texttt{PatsIota} is preserved by any environment extension that leaves \texttt{pats}
unchanged and only grows \texttt{constants}.
\thesisref{no direct counterpart; proof infrastructure}
\end{lemma}
\begin{proof}
\leanok
Rewrite each \texttt{pats} hypothesis by the equation and transport the constant lookups
along the environment order.
\end{proof}

\begin{theorem}[\texttt{PatsIota} survives \texttt{addInduct}]
\label{thm:indparams-patsiota-induct}
\lean{Lean4Lean.VEnv.PatsIota.induct}
\leanok
\contrib{iota} \srcloc{Lean4Lean/Theory/Typing/InductiveParams.lean}{181}
\uses{structure:indparams-patsiota, def:ind-add-induct, structure:ind-wf}
If $env$ satisfies \texttt{PatsIota} and $decl$ is well-formed in $env$, then the result of
$env.\mathsf{addInduct}\;decl$ satisfies \texttt{PatsIota}.
\thesisref{no direct counterpart; the inductive step of the registry invariant}
\end{theorem}
\begin{proof}
\uses{family:indlem-pats-origin, family:indlem-addinduct-consts, thm:indlem-rules-ctor-shape, thm:indlem-iota-topattern-inj, family:ind-shape-corollaries, family:indlem-addconst-spec}
\leanok
By \texttt{addInduct\_pats\_origin} each entry is old or new. A new entry is keyed by a
recursor of the block, whose type is \texttt{RecHeaded} by \texttt{rec\_shape.recHeaded} and
whose constructor is registered with \texttt{CtorShape} by \texttt{addInduct\_rule\_ctor}; an
old entry's recursor cannot be a new one because \texttt{addInduct\_rec\_fresh} says the new
names were unbound; and two new entries with the same key come from the same recursor, by
\texttt{addInduct\_recs\_name\_inj}, and from the same rule, by \texttt{rules\_nodup} with
\texttt{nodup\_map\_inj\_on}, so their reducts coincide. At about fifty lines this is the
most substantial proof in the group and the one that actually consumes the
\texttt{VInductDecl.WF} clauses.
\end{proof}

\begin{theorem}[Every well-formed environment satisfies the registry invariant]
\label{thm:indparams-wf-patsiota}
\lean{Lean4Lean.VEnv.WF.patsIota}
\leanok
\contrib{iota} \srcloc{Lean4Lean/Theory/Typing/InductiveParams.lean}{232}
\uses{structure:indparams-patsiota, def:venv-wf}
$env.\mathsf{WF}$ implies $env.\mathsf{PatsIota}$.
\thesisref{no direct counterpart; the registry invariant, globally}
\end{theorem}
\begin{proof}
\uses{thm:indparams-patsiota-induct, thm:indparams-patsiota-of-le, family:indparams-pats-eq-or-induct}
\leanok
Induction on the \texttt{WF'} chain: the empty environment has no registered patterns, and
each step is either \texttt{PatsIota.of\_le} or \texttt{PatsIota.induct}, according to
\texttt{pats\_eq\_or\_induct}.
\end{proof}

\begin{lemma}[The five structural Params side conditions, discharged]
\label{family:indparams-params-sides}
\lean{Lean4Lean.VEnv.WF.pat_simple, Lean4Lean.VEnv.WF.pat_uniq, Lean4Lean.VEnv.WF.pat_app_l, Lean4Lean.VEnv.WF.pat_app_l_uniq, Lean4Lean.VEnv.WF.pat_app_uniq}
\leanok
\contrib{iota} \srcloc{Lean4Lean/Theory/Typing/InductiveParams.lean}{249}
\uses{class:cr-params, def:venv-wf, inductive:pattern-core}
For a well-formed environment: every registered pattern is a \texttt{SimplePattern}; if a
subpattern of one registered redex intersects another registered redex then the two redexes
and the subpattern all coincide and the reducts agree; the recursor spine of a redex has no
application subpattern of its own; a variable slot of one recursor spine never intersects
another recursor spine; and a subpattern of one recursor spine never intersects a subpattern
of another constructor spine.
\thesisref{no direct counterpart; the coherence conditions of the abstract development}
\end{lemma}
\begin{proof}
\uses{thm:indparams-wf-patsiota, thm:indparams-patsiota-rec-ne-ctor, family:pattern-spine-combinatorics, family:indparams-inter, family:indparams-iota-subpattern}
\leanok
All five reduce through \texttt{PatsIota.shape} to the spine combinatorics.
\texttt{pat\_uniq} splits on where the subpattern sits: at the top, \texttt{functional}
closes; inside the recursor spine, the arity would drop below the one \texttt{arity} fixes;
inside the constructor spine, a recursor would be a constructor. \texttt{pat\_app\_l\_uniq}
and \texttt{pat\_app\_uniq} are \texttt{varN\_const\_inter} plus \texttt{arity},
respectively \texttt{rec\_ne\_ctor}. These are the real content of the instance and they
need no extra hypotheses.
\end{proof}

\begin{definition}[Every definitional axiom is realised by a registered pattern, \texttt{DefEqsAsPats}]
\label{def:indparams-defeqs-as-pats}
\lean{Lean4Lean.VEnv.DefEqsAsPats}
\leanok
\contrib{iota} \srcloc{Lean4Lean/Theory/Typing/InductiveParams.lean}{378}
\uses{inductive:pattern-matches, inductive:pattern-rhs, def:hastype-istype, class:cr-params}
$env.\mathsf{DefEqsAsPats}\;U$ asserts, in the shape of the \texttt{Params} field
\texttt{extra\_pat}, that for every $env.\mathsf{defeqs}\;df$ and every well-formed level
instantiation there is a registered pattern matching the instantiated left-hand side whose
side conditions hold and whose reduct is the instantiated right-hand side.
\thesisref{no direct counterpart; the gap between \texttt{defeqs} and \texttt{pats}}
\reviewnote{The 27-line docstring (\srcloc{Lean4Lean/Theory/Typing/InductiveParams.lean}{351})
is the most useful documentation in the group. It explains that $\delta$ rules and the
quotient rule live in \texttt{defeqs}, why registering \texttt{SimplePattern.defn} would
require a new environment invariant separating definition names from recursor and
constructor names, and why the quotient redex is not a \texttt{SimplePattern} at all. The
consequence, stated plainly, is that \texttt{DefEqsAsPats} fails for any environment
containing a single \texttt{def} or \texttt{quot}. Two thirds of it (lines 355, 362--377) is
tagged \texttt{trproj} by the blame and is, unlike every other \texttt{trproj}-tagged block
in this group, genuinely absent from \texttt{iota}, where the hypothesis is called a
``design hypothesis'' and the $\delta$ analysis is missing.}
\end{definition}

\begin{definition}[The Params instance induced by an environment, \texttt{toParams}]
\label{def:indparams-toparams}
\lean{Lean4Lean.VEnv.toParams}
\leanok
\contrib{iota} \srcloc{Lean4Lean/Theory/Typing/InductiveParams.lean}{393}
\uses{family:indparams-params-sides, thm:pattern-ok-exists-realizer, def:indparams-defeqs-as-pats, class:cr-params, rule:isdefeq-pat}
For a well-formed $env$ satisfying $\mathsf{DefEqsAsPats}\;U$, $env.\mathsf{toParams}$ is a
\texttt{Params} structure whose abstract reduction relation is $env.\mathsf{pats}$: the five
structural conditions come from \ref{family:indparams-params-sides}, \texttt{pat\_wf} is
\texttt{IsDefEq.pat} applied to a realizer extracted by
\ref{thm:pattern-ok-exists-realizer}, \texttt{pat\_env} is the identity, and
\texttt{extra\_pat} is the hypothesis.
\thesisref{no direct counterpart; the bridge to the Church-Rosser development}
\reviewnote{This is the point of the file: the abstract development had no instance at all
before. Its value is limited by \texttt{DefEqsAsPats}, which as it stands restricts
\texttt{toParams} to environments with no definitional axiom whatsoever.}
\end{definition}

\begin{lemma}[\texttt{DefEqsAsPats} holds vacuously without definitional axioms]
\label{thm:indparams-defeqs-as-pats-vacuous}
\lean{Lean4Lean.VEnv.DefEqsAsPats.of_no_defeqs}
\leanok
\contrib{trproj} \srcloc{Lean4Lean/Theory/Typing/InductiveParams.lean}{413}
\uses{def:indparams-defeqs-as-pats}
An environment whose \texttt{defeqs} relation is empty satisfies
$\mathsf{DefEqsAsPats}\;U$ for every $U$.
\thesisref{no direct counterpart}
\end{lemma}
\begin{proof}
\leanok
One line: \texttt{absurd} on the \texttt{defeqs} hypothesis. The line blame tags this
\texttt{trproj}; identical code is on \texttt{iota}.
\end{proof}

\begin{definition}[A Params instance for a single inductive block, \texttt{inductParams}]
\label{def:indparams-induct-params}
\lean{Lean4Lean.VEnv.inductParams}
\leanok
\contrib{trproj} \srcloc{Lean4Lean/Theory/Typing/InductiveParams.lean}{419}
\uses{def:indparams-toparams, thm:indparams-defeqs-as-pats-vacuous, family:indlem-pats-defeqs-preserved, def:venv-wf}
For a declaration well-formed in the empty environment with
$\mathsf{addInduct}\;\emptyset\;decl = \mathsf{some}\;env$, \texttt{inductParams} is the
\texttt{Params} instance at $env$: well-formedness is witnessed by the one-declaration chain,
and \texttt{DefEqsAsPats} holds vacuously because \texttt{addInduct} adds no definitional
axiom.
\thesisref{no direct counterpart; an inhabitation witness}
\reviewnote{The only witness that \texttt{Params} is inhabitable at all, which is genuinely
useful as a sanity check, but the environment is a toy: one inductive block and nothing
else. The line blame tags it \texttt{trproj}; identical on \texttt{iota}.}
\end{definition}

\begin{lemma}[Church-Rosser for a one-block environment]
\label{thm:indparams-cr-defeq-of-induct}
\lean{Lean4Lean.VEnv.IsDefEq.crDefEq_of_induct}
\leanok
\contrib{trproj} \srcloc{Lean4Lean/Theory/Typing/InductiveParams.lean}{426}
\uses{def:indparams-induct-params, thm:isdefeq-church-rosser, def:crdefeq, def:normaleq, def:onctx}
In an environment obtained by adding one well-formed inductive block to the empty
environment, definitionally equal terms in a well-formed context are related by
\texttt{CRDefEq}: they have parallel reduction sequences meeting at \texttt{NormalEq} terms.
\thesisref{no direct counterpart; a demonstration that the instance can be used}
\end{lemma}
\begin{proof}
\leanok
Instantiate \texttt{inductParams} and apply \texttt{IsDefEq.church\_rosser}.
\axfoot{sorryAx via VEnv.NormalEq.parRed (ChurchRosser.lean:1193 and :1212, master);
Classical.choice; Quot.sound; propext}
\reviewnote{The statement reads as a finished result and is not one.
\texttt{IsDefEq.church\_rosser} routes through \texttt{NormalEq.parRed}, which still contains
two \texttt{sorry}s, so the conclusion is conditional on them; the docstring says nothing
about that. Nothing consumes the theorem.}
\end{proof}

\section{Contribution summary}

The $\iota$ contribution to this group does four things. First, it replaces two
\texttt{sorry}-definitions by a real specification: \ref{def:ind-add-induct} stages the
environment extension into type formers, constructors, recursors and rules, and
\ref{structure:ind-wf} states in twenty clauses what the kernel checks, each clause typed in
the environment the kernel checks it in and each documented with its thesis reference. The
syntactic half of that record is \ref{def:ind-ctor-result} to \ref{def:ind-rule-shape}, whose
de Bruijn arithmetic is correct and is exercised against real kernel declarations by
\texttt{Tests/IotaShape.lean}. Second, it makes the $\iota$ rule a first-class reduction rule:
\ref{def:pattern-rhs-generic} turns the thesis's ``all substitution instances'' parenthesis
into one typing obligation, \ref{def:pattern-check-realizes} makes the side conditions a
positive predicate so that the new \texttt{IsDefEq.pat} constructor is legal,
\ref{inductive:pattern-template-headed} records that a registered rule computes, and
\ref{def:pattern-iota-rhs} with \ref{thm:pattern-iota-rhs-apply} identify the model's reduct
with the kernel's \texttt{inductiveReduceRecCore} reduct argument slice for argument slice.
Third, \ref{thm:indlem-add-induct-wf} proves that adding a well-formed block preserves
orderedness, with \ref{family:indlem-pat-registration} and \ref{family:indlem-pats-origin}
giving both directions of ``\texttt{pats} holds exactly the block's $\iota$ rules''. Fourth,
\ref{structure:indparams-patsiota} to \ref{def:indparams-toparams} build the first concrete
instance of the parameter record the Church-Rosser development is stated over.

What remains open. Nothing derives \ref{structure:ind-wf} from
\texttt{Lean4Lean/Inductive/Add.lean}, the repository's own re-implementation of the kernel's
inductive checker, so \texttt{TrEnv'.induct} still takes it as a hypothesis; two of the
model's deliberate deviations, manifest binders in positivity and the exclusion of nested
inductives, actively obstruct closing that gap. \ref{def:indparams-toparams} applies only to
environments with no definitional axiom, so the only instance built,
\ref{def:indparams-induct-params}, is a toy. \ref{thm:indparams-wf-pats-origin} is groundwork
for \texttt{VEnv.WF.patsStrong}, which is still \texttt{sorry}. And K-like reduction is not
modelled at all, so for subsingleton eliminators the model's reduction relation is strictly
weaker than the kernel's.

\section{Review notes}
\label{sec:ind-review}

\begin{itemize}
\item \textbf{High.} \srcloc{Lean4Lean/Verify/Environment/Basic.lean}{583}: nothing derives
\texttt{VInductDecl.WF} from the kernel's own inductive checker. \texttt{TrEnv'.induct} takes
\texttt{decl.WF env} as a hypothesis and \texttt{Lean4Lean/Inductive/Add.lean} is never
related to it, so the refinement chain still has an assumption at the hardest declaration
kind. This is better than \texttt{master}, where the predicate was \texttt{sorry} and the
assumption was uninformative, and \texttt{Tests/IotaShape.lean} validates the record by
\texttt{decide} on concrete kernel data, but that is evidence, not a proof. The contribution
should be described as specified and validated, not verified.
\item \textbf{High.} \srcloc{Lean4Lean/Theory/Typing/InductiveParams.lean}{426}:
\texttt{IsDefEq.crDefEq\_of\_induct} is presented as Church-Rosser for such an environment,
but \texttt{IsDefEq.church\_rosser} depends on the two \texttt{sorry}s in
\texttt{NormalEq.parRed} (\srcloc{Lean4Lean/Theory/Typing/ChurchRosser.lean}{1193} and
\srcloc{Lean4Lean/Theory/Typing/ChurchRosser.lean}{1212}). The docstring gives no hint that
the conclusion is conditional, and no other declaration uses the theorem.
\item \textbf{Medium.} \srcloc{Lean4Lean/Theory/Typing/InductiveParams.lean}{393}:
\texttt{toParams} is a \texttt{Params} instance only for environments satisfying
\texttt{DefEqsAsPats}, that is, containing no \texttt{def}, \texttt{mutualDef} or
\texttt{quot} at all. Every realistic Lean environment fails it, so the Church-Rosser
development is still not connected to the model's own well-formed environments. Documented at
length in the \texttt{DefEqsAsPats} docstring, but it is the main limitation of the file.
\item \textbf{Medium.} \srcloc{Lean4Lean/Theory/Inductive.lean}{174}:
\texttt{VExpr.RuleShape} pins only the number of arguments the reduct passes to the minor
premise after the fields, not their shape, where the thesis requires
$v_i = \lambda x{::}\xi_i.\;\mathrm{rec}_P\;C\;e\;\pi_i[b,x]\;(u_i\;x)$. Together with
\texttt{rules\_wf}, which only forces them to be well-typed, \texttt{VInductDecl.WF} accepts
$\iota$ rules the kernel would never generate. Flagged in the docstring, but it caps what any
soundness theorem proved against this model can mean.
\item \textbf{Medium.} \srcloc{Lean4Lean/Theory/Inductive.lean}{257}: K-like reduction is not
modelled. \texttt{addRecRule} registers only the constructor rule and \texttt{VRecursor.k} is
recorded but never used, while the kernel does perform it (\texttt{toCtorWhenK},
\srcloc{Lean4Lean/Inductive/Reduce.lean}{29}, called from
\srcloc{Lean4Lean/Inductive/Reduce.lean}{110}). For K-like recursors the model's reduction
relation is strictly weaker than the kernel's and the verified WHNF cannot be shown complete
for them.
\item \textbf{Medium.} \srcloc{Lean4Lean/Theory/Inductive.lean}{93}:
\texttt{FieldPositive}, \texttt{CtorPositive} and \texttt{WF.universes} read the manifest
$\Pi$-binders of a field type, whereas the kernel reduces to weak head normal form at each
step of \texttt{checkPositivity} (\srcloc{Lean4Lean/Inductive/Add.lean}{189}). The model is
therefore strictly stricter than the kernel: a declaration whose field type is a reducible
definition unfolding to a $\Pi$ passes the kernel and fails \texttt{VInductDecl.WF}.
Documented, but it is an obstacle to ever discharging the first item above.
\item \textbf{Medium.} \srcloc{Lean4Lean/Theory/Inductive.lean}{335}:
\texttt{recs\_over\_block} and \texttt{rules\_ctor} require every recursor to eliminate one of
the block's own type formers, which excludes nested inductives. \texttt{Environment.addInductive}
does add such blocks, so \texttt{TrEnv'} can never be constructed for an environment containing
one. Documented as future work, namely modelling the kernel's \texttt{ElimNestedInductive} pass.
\item \textbf{Low.} \srcloc{Lean4Lean/Theory/Inductive.lean}{196}: dead code.
\texttt{VExpr.RecShape.one\_le\_numMotives} (line 196),
\texttt{VExpr.RecShape.majorFormer?\_eq} (line 202) and
\texttt{VInductDecl.LargeElim.shape} (line 242) are referenced nowhere in the repository, and
so is the whole \texttt{Pattern.RHS.spine} and \texttt{iotaCounts} cluster at
\srcloc{Lean4Lean/Theory/Typing/Pattern.lean}{509} to 678, eight declarations and about
seventy lines used only by each other.
\item \textbf{Low.} \srcloc{Lean4Lean/Theory/Typing/InductiveLemmas.lean}{451}: placement and
naming. \texttt{nodup\_map\_inj\_on} is a pure \texttt{List} lemma declared as
\texttt{Lean4Lean.VEnv.nodup\_map\_inj\_on}; the generic \texttt{foldlM} and
\texttt{addConst\_foldlM} toolkit and the \texttt{addQuot} and \texttt{addDefEqs} lemmas,
which are about quotients and definitions, all live in a file named \texttt{InductiveLemmas};
and \texttt{VEnv.addRecRule} (\srcloc{Lean4Lean/Theory/Inductive.lean}{259}) spells the major
index out instead of using \texttt{r.getMajorIdx}, which every lemma about it does use.
Similarly \texttt{Pattern.inter\_app\_const} and its two siblings are \texttt{\_root\_}
\texttt{Pattern} lemmas declared inside \texttt{namespace VEnv}.
\item \textbf{Low.} \srcloc{Lean4Lean/Theory/Typing/Pattern.lean}{283}: redundancy.
\texttt{Pattern.Matches.uniq} (line 143) and \texttt{Pattern.matches\_determ} (line 283) are
the same statement proved twice; both are pre-existing \texttt{master} code that the branches
left untouched (\texttt{git diff master} on this file is 493 insertions and no deletions), so
the duplication is inherited rather than introduced.
Similarly \texttt{MotiveShape}'s ``ends in a sort'' clause is
subsumed by \texttt{VInductDecl.WF.recs\_elim}, which pins the exact sort, and
\texttt{WF.universes} bundles the purely syntactic bound
\texttt{decl.nparams} $\le$ \texttt{t.type.piArity} into a clause guarded by stage 0
succeeding.
\item \textbf{Low, fidelity.} \srcloc{Lean4Lean/Theory/Inductive.lean}{75}:
\texttt{ValidIndApp} drops the kernel's exact-arity check on the occurrence's arguments,
\texttt{args.size == params.size + nindices[i]}
(\srcloc{Lean4Lean/Inductive/Add.lean}{159}), although its docstring claims to mirror
\texttt{isValidIndApp?}; \texttt{FieldInIndices}
(\srcloc{Lean4Lean/Theory/Inductive.lean}{107}) searches only the arguments past
\texttt{nparams} where the kernel searches all of them
(\srcloc{Lean4Lean/Inductive/Add.lean}{275}), the two agreeing only by an argument about
\texttt{CtorResult} that is recorded nowhere; and \texttt{LargeElim} demands a typing at
$\mathsf{sort}\;0$ where the kernel tests \texttt{isAlwaysZero} on the inferred level
(\srcloc{Lean4Lean/Inductive/Add.lean}{271}), which agree only up to level defeq and
uniqueness of typing.
\item \textbf{Low, undocumented limitations.}
\srcloc{Lean4Lean/Theory/Inductive.lean}{351}: \texttt{WF.universes} hides the purely
syntactic bound \texttt{decl.nparams} $\le$ \texttt{t.type.piArity} inside a clause guarded
by \texttt{addTypes} succeeding, so the arity fact is unavailable when stage 0 fails; and no
clause of \texttt{VInductDecl.WF} mentions \texttt{VRecursor.all}, which the model's own
well-formedness therefore never uses (the refinement pins it to the kernel's separately,
\srcloc{Lean4Lean/Verify/Environment/Basic.lean}{240}). Unlike the nested-inductive, K-like,
$\eta$, \texttt{whnf} and \texttt{rule\_shape} limitations, neither is recorded in the
source.
\end{itemize}
